\setrunningtitles{Januarius}{January}
\newMonth{Januarius}

% Set initial running titles for left/right pages:


% ---- martyrology/mart01/mart0101.htm
\needspace{10\baselineskip}
\begin{paracol}{2}
\selectlanguage{latin}
\begin{center}{\color{gregoriocolor} Kaléndis Januárii. Luna\dots\ }\end{center}
\switchcolumn
\selectlanguage{english}
\begin{center}{\color{gregoriocolor} The  1\textsuperscript{st} Day of 
 January. The\dots\ Day of the Moon.}\end{center}
\end{paracol}

% TODO: find out a way we can distribute the letter-number pairs evenly.
\noindent\begin{tabularx}{\linewidth}{*{19}{>{\centering\arraybackslash}X}}
 \textcolor{gregoriocolor}{a} & \textcolor{gregoriocolor}{b} & \textcolor{gregoriocolor}{c} & \textcolor{gregoriocolor}{d} & \textcolor{gregoriocolor}{e} & \textcolor{gregoriocolor}{f} & \textcolor{gregoriocolor}{g} & \textcolor{gregoriocolor}{h} & \textcolor{gregoriocolor}{i} & \textcolor{gregoriocolor}{k} & \textcolor{gregoriocolor}{l} & \textcolor{gregoriocolor}{m} & \textcolor{gregoriocolor}{n} & \textcolor{gregoriocolor}{p} & \textcolor{gregoriocolor}{q} & \textcolor{gregoriocolor}{r} & \textcolor{gregoriocolor}{s} & \textcolor{gregoriocolor}{t} & \textcolor{gregoriocolor}{u} \\
 2 & 3 & 4 & 5 & 6 & 7 & 8 & 9 & 10 & 11 & 12 & 13 & 14 & 15 & 16 & 17 & 18 & 19 & 20 \\
\end{tabularx}
\vspace{0.5\baselineskip}
\noindent\begin{tabularx}{\linewidth}{*{12}{>{\centering\arraybackslash}X}}
 \textcolor{gregoriocolor}{A} & \textcolor{gregoriocolor}{B} & \textcolor{gregoriocolor}{C} & \textcolor{gregoriocolor}{D} & \textcolor{gregoriocolor}{E} & F & \textcolor{gregoriocolor}{F} & \textcolor{gregoriocolor}{G} & \textcolor{gregoriocolor}{H} & \textcolor{gregoriocolor}{M} & \textcolor{gregoriocolor}{N} & \textcolor{gregoriocolor}{P} \\
 21 & 22 & 23 & 24 & 25 & 26 & 26 & 27 & 28 & 29 & 30 & 1 \\
\end{tabularx}

\begin{paracol}{2}
\selectlanguage{latin}
\lettrine[lines=2]{C}{ircumcísio} Dómini nostri Jesu Christi, et Octáva Nativitátis ejúsdem.
\switchcolumn
\selectlanguage{english}
\lettrine[lines=2]{T}{he} Circumcision of our Lord Jesus Christ, and the octave of his Nativity.
\switchcolumn*
\selectlanguage{latin}
Romæ pássio sanctæ Martínæ, Vírginis et 
 Mártyris; quæ, sub Alexándro Imperatóre, divérsis tormentórum genéribus 
 cruciáta, tandem, gládio percússa, martyrii palmam adépta est. Ipsíus 
 vero festum tértio Kaléndas Februárii recólitur.
\switchcolumn
\selectlanguage{english}
At Rome, under Emperor Alexander, St. Martina, virgin, who endured various 
 kinds of torments, and being beheaded, received the palm of martyrdom. 
 Her feast is kept on the 30th of this month.
\switchcolumn*
\selectlanguage{latin}
Cæsaréæ, in Cappadócia, deposítio sancti Basilíi, cognoménto Magni, Epíscopi, Confessóris et Ecclésiæ Doctóris; qui, témpore Valéntis Imperatóris, doctrína et sapiéntia insignítus omnibúsque virtútibus exornátus, mirabíliter effúlsit, et Ecclésiam advérsus Ariános et Macedoniános inexpugnábili constántia deféndit. Ejus autem festívitas potíssimum ágitur décimo octávo Kaléndas Júlii, quo die Epíscopus ordinátus est.
\switchcolumn
\selectlanguage{english}
At Caesarea in Cappadocia, the death of St. Basil the Great, bishop, 
 confessor, and doctor of the Church, renowned for his learning and wisdom 
 and gifted with every virtue, who during the reign of Emperor Valens 
 wonderfully displayed his talents as he defended the Church with great 
 constancy against the Arians and Macedonians. His feast, however, is 
 appropriately kept on the 14th of June, the day on which he was consecrated 
 bishop.
\switchcolumn*
\selectlanguage{latin}
Apud montem Senárium, in Etrúria, natális sancti Bonfílii Confessóris, e 
 septem Fundatóribus Ordinis Servórum beátæ 
 Maríæ Vírginis, quam cum diem impénse coluísset, ab ipsa in cælum repénte 
 evocátus est. Illíus porro ac Sociórum festum prídie Idus Februárii 
 celebrátur.
\switchcolumn
\selectlanguage{english}
In Tuscany, on Mount Senario, St. Bonfilius, confessor, one of the seven 
 founders of the Order of the Servites of the Blessed Virgin Mary, who, 
 having honoured her devoutly, was suddenly called to heaven by her. 
 His feast, with that of his companions, is kept on February 12th.
\switchcolumn*
\selectlanguage{latin}
Romæ sancti Almáchii Mártyris, qui, cum díceret: 
 « Hódie Octávæ Domínici diéi sunt, cessáte a superstitiónibus idolórum et a 
 sacrifíciis pollútis », proptérea, jubénte Præfécto Urbis Alípio, a 
 gladiatóribus occísus est.
\switchcolumn
\selectlanguage{english}
At Rome, St. Almachius, martyr, who, by the command of Alipius, governor of 
 the city, was killed by the gladiators for saying, ``Today is the Octave of 
 our Lord's birth; put an end to the worship of idols, and abstain from 
 unclean sacrifices.''
\switchcolumn*
\selectlanguage{latin}
Item Romæ, via Appia, corónæ sanctórum mílitum 
 trigínta Mártyrum, sub Diocletiáno Imperatóre.
\switchcolumn
\selectlanguage{english}
In the same city, on the Appian Way, the crowning with martyrdom of thirty 
 holy soldiers under Emperor Diocletian.
\switchcolumn*
\selectlanguage{latin}
Apud Spolétum sancti Concórdii, Presbyteri et Mártyris; qui, tempóribus 
 Antoníni Imperatóris, primo cæsus fústibus, 
 dehinc equúleo suspénsus, ac póstea macerátus in cárcere, ibíque Angélica 
 visitatióne confortátus, demum gládio vitam finívit.
\switchcolumn
\selectlanguage{english}
At Spoleto, in the time of Emperor Antoninus, St. Concordius, priest and 
 martyr, who was beaten with clubs, then stretched on the rack, and after a 
 long confinement in prison, where he was visited by an angel, lost his life 
 by the sword.
\switchcolumn*
\selectlanguage{latin}
Eódem die sancti Magni Mártyris.
\switchcolumn
\selectlanguage{english}
The same day, St. Magnus, martyr.
\switchcolumn*
\selectlanguage{latin}
In Africa beáti Fulgéntii, Ruspénsis Ecclésiæ 
 Epíscopi, qui, témpore Wandálicæ persecutiónis, ob cathólicam fidem 
 eximiámque doctrínam, ab Ariánis multa perpéssus et in Sardíniam relegátus 
 est; atque tandem, ad própriam Ecclésiam redíre permíssus, vita et verbo 
 clarus, sancto fine quiévit.
\switchcolumn
\selectlanguage{english}
In Africa, St. Fulgentius, bishop of Rusp, who suffered much from the 
 Arians, during the persecution of the Vandals, for holding the Catholic 
 faith and teaching an excellent doctrine. After being banished to 
 Sardinia, he was permitted to return to his diocese, where he ended his life 
 by a holy death, leaving a reputation for sanctity and eloquence.
\switchcolumn*
\selectlanguage{latin}
Teáte, in Aprútio citerióre, natális sancti Justíni, ejúsdem civitátis 
 Epíscopi, sanctitáte vitæ ac miráculis clari.
\switchcolumn
\selectlanguage{english}
At Chieti in Abruzzo, the birthday of St. Justin, bishop of that city, 
 illustrious for holiness of life and for his miracles.
\switchcolumn*
\selectlanguage{latin}
In território Lugdunénsi, monastério Jurénsium, sancti Eugéndi Abbátis, 
 cujus vita virtútibus et miráculis plena refúlsit.
\switchcolumn
\selectlanguage{english}
In the diocese of Lyons, in the monastery of St. Claude, St. Eugendus, 
 abbot, whose life was eminent for virtues and miracles.
\switchcolumn*
\selectlanguage{latin}
Apud Silviníacum, in Gállia, sancti Odilónis, 
 Abbátis Cluniacénsis, qui primus Commemoratiónem ómnium Fidélium Defunctórum, 
 prima die post festum ómnium Sanctórum, in suis monastériis fíeri præcépit; 
 quem ritum póstea universális Ecclésia recípiens comprobávit.
\switchcolumn
\selectlanguage{english}
At Souvigny in France, St. Odilo, abbot of Cluny, who was the first to 
 prescribe that the commemoration of all the faithful departed should be made 
 in his monasteries the day after the feast of All Saints. This 
 practice was afterwards received and approved by the universal Church.
\switchcolumn*
\selectlanguage{latin}
Romæ natális sancti Vincéntii Maríæ Strambi, 
 Epíscopi Maceraténsis et Tolentíni, Congregatiónis a Cruce et Passióne Jesu 
 Sodális, pastoráli zelo præclári, quem Pius Papa Duodécimus inter Sanctos 
 rétulit.
\switchcolumn
\selectlanguage{english}
At Rome, the birthday of St. Vincent Maria Strambi, Bishop of Macerata and 
 Tolentino, of the Order of Passionists, renowned for his pastoral zeal, whom 
 Pope Pius XII numbered among the saints.
\switchcolumn*
\selectlanguage{latin}
Alexandríæ deposítio sanctæ Euphrósynæ Vírginis, 
 quæ in monastério virtúte abstinéntiæ ac miráculis cláruit.
\switchcolumn
\selectlanguage{english}
At Alexandria, the departure from this world of St. Euphrosyna, virgin, who 
 was renowned in her monastery for the virtue of abstinence, and for the gift 
 of miracles.
\switchcolumn*
\selectlanguage{latin}
Et álibi aliórum 
 plurimórum sanctórum Mártyrum et Confessórum, atque sanctárum Vírginum.

\textcolor{gregoriocolor}{\Rbar.}~Deo grátias.
\switchcolumn
\selectlanguage{english}
And elsewhere in divers 
 places, many other holy martyrs, confessors, and holy virgins.

\textcolor{gregoriocolor}{\Rbar.}~Thanks be to God.
\switchcolumn*
\selectlanguage{latin}
\end{paracol}

% \continueMonth{Januarius}


% ---- martyrology/mart01/mart0102.htm
\needspace{10\baselineskip}
\begin{paracol}{2}
\selectlanguage{latin}
\begin{center}{\color{gregoriocolor} Quarto Nonas Januárii. 
 Luna\dots\ }\end{center}
\switchcolumn
\selectlanguage{english}
\begin{center}{\color{gregoriocolor} The   2\textsuperscript{nd} Day of 
 January. The\dots\ Day of the Moon.}\end{center}
\end{paracol}

\noindent\begin{tabularx}{\linewidth}{*{19}{>{\centering\arraybackslash}X}}
 \textcolor{gregoriocolor}{a} & \textcolor{gregoriocolor}{b} & \textcolor{gregoriocolor}{c} & \textcolor{gregoriocolor}{d} & \textcolor{gregoriocolor}{e} & \textcolor{gregoriocolor}{f} & \textcolor{gregoriocolor}{g} & \textcolor{gregoriocolor}{h} & \textcolor{gregoriocolor}{i} & \textcolor{gregoriocolor}{k} & \textcolor{gregoriocolor}{l} & \textcolor{gregoriocolor}{m} & \textcolor{gregoriocolor}{n} & \textcolor{gregoriocolor}{p} & \textcolor{gregoriocolor}{q} & \textcolor{gregoriocolor}{r} & \textcolor{gregoriocolor}{s} & \textcolor{gregoriocolor}{t} & \textcolor{gregoriocolor}{u} \\
 3 & 4 & 5 & 6 & 7 & 8 & 9 & 10 & 11 & 12 & 13 & 14 & 15 & 16 & 17 & 18 & 19 & 20 & 21 \\
\end{tabularx}
\vspace{0.5\baselineskip}
\noindent\begin{tabularx}{\linewidth}{*{12}{>{\centering\arraybackslash}X}}
 \textcolor{gregoriocolor}{A} & \textcolor{gregoriocolor}{B} & \textcolor{gregoriocolor}{C} & \textcolor{gregoriocolor}{D} & \textcolor{gregoriocolor}{E} & F & \textcolor{gregoriocolor}{F} & \textcolor{gregoriocolor}{G} & \textcolor{gregoriocolor}{H} & \textcolor{gregoriocolor}{M} & \textcolor{gregoriocolor}{N} & \textcolor{gregoriocolor}{P} \\
 22 & 23 & 24 & 25 & 26 & 27 & 27 & 28 & 29 & 30 & 1 & 2 \\
\end{tabularx}

\begin{paracol}{2}
\selectlanguage{latin}
\lettrine[lines=1]{O}{ctáva} sancti Stéphani Protomártyris.
\switchcolumn
\selectlanguage{english}
\lettrine[lines=1]{T}{he} Octave of St. Stephen, the first martyr.
\switchcolumn*
\selectlanguage{latin}
Romæ commemorátio plurimórum sanctórum Mártyrum, 
 qui, spreto Diocletiáni Imperatóris edícto quo tradi sacri Códices 
 jubebántur, pótius córpora carnifícibus quam sancta dare cánibus maluérunt.
\switchcolumn
\selectlanguage{english}
At Rome, the commemoration of many holy martyrs, who, despising the edict of 
 Emperor Diocletian, which ordered that the sacred books should be delivered 
 up, preferred to offer their bodies to the executioners rather than to give 
 holy things to dogs.
\switchcolumn*
\selectlanguage{latin}
Antiochíæ pássio beáti Isidóri Epíscopi.
\switchcolumn
\selectlanguage{english}
At Antioch, the passion of blessed Isidore, bishop.
\switchcolumn*
\selectlanguage{latin}
Tomis, in Ponto, sanctórum fratrum Argéi, 
 Narcíssi et Marcellíni púeri. Hic, sub Licínio Príncipe, cum inter 
 tirónes esset comprehénsus et nollet militáre, hinc, cæsus ad mortem ac diu 
 macerátus in cárcere, demum, in mare demérsus, martyrium consummávit; ejus 
 autem fratres gládio perémpti sunt.
\switchcolumn
\selectlanguage{english}
At Tomis in Pontus, in the time of Emperor Licinius, three holy brothers, 
 Argeus, Narcissus, and the young man Marcellinus. This last, being 
 enrolled among the new soldiers, and refusing to serve, was beaten almost to 
 death, and for a long while kept in prison. Being finally cast into 
 the sea, he finished his martyrdom, and his brothers were beheaded.
\switchcolumn*
\selectlanguage{latin}
Medioláni sancti Martiniáni Epíscopi.
\switchcolumn
\selectlanguage{english}
At Milan, St. Martinian, bishop.
\switchcolumn*
\selectlanguage{latin}
Nítriæ, in Ægypto, beáti Isidóri, Epíscopi et 
 Confessóris.
\switchcolumn
\selectlanguage{english}
In Nitria in Egypt, blessed Isidore, bishop and confessor.
\switchcolumn*
\selectlanguage{latin}
Ipso die sancti Siridiónis Epíscopi.
\switchcolumn
\selectlanguage{english}
The same day, St. Siridion, bishop.
\switchcolumn*
\selectlanguage{latin}
In Thebáide sancti Macárii Alexandríni, Presbyteri et Abbátis.
\switchcolumn
\selectlanguage{english}
In Thebais, St. Macarius of Alexandria, abbot.
\switchcolumn*
\selectlanguage{latin}
\end{paracol}


% ---- martyrology/mart01/mart0103.htm
\needspace{10\baselineskip}
\begin{paracol}{2}
\selectlanguage{latin}
\begin{center}{\color{gregoriocolor} Tértio Nonas Januárii. 
 Luna\dots\ }\end{center}
\switchcolumn
\selectlanguage{english}
\begin{center}{\color{gregoriocolor} The   3\textsuperscript{rd} Day of 
 January. The\dots\ Day of the Moon.}\end{center}
\end{paracol}

\noindent\begin{tabularx}{\linewidth}{*{19}{>{\centering\arraybackslash}X}}
 \textcolor{gregoriocolor}{a} & \textcolor{gregoriocolor}{b} & \textcolor{gregoriocolor}{c} & \textcolor{gregoriocolor}{d} & \textcolor{gregoriocolor}{e} & \textcolor{gregoriocolor}{f} & \textcolor{gregoriocolor}{g} & \textcolor{gregoriocolor}{h} & \textcolor{gregoriocolor}{i} & \textcolor{gregoriocolor}{k} & \textcolor{gregoriocolor}{l} & \textcolor{gregoriocolor}{m} & \textcolor{gregoriocolor}{n} & \textcolor{gregoriocolor}{p} & \textcolor{gregoriocolor}{q} & \textcolor{gregoriocolor}{r} & \textcolor{gregoriocolor}{s} & \textcolor{gregoriocolor}{t} & \textcolor{gregoriocolor}{u} \\
 4 & 5 & 6 & 7 & 8 & 9 & 10 & 11 & 12 & 13 & 14 & 15 & 16 & 17 & 18 & 19 & 20 & 21 & 22 \\
\end{tabularx}
\vspace{0.5\baselineskip}
\noindent\begin{tabularx}{\linewidth}{*{12}{>{\centering\arraybackslash}X}}
 \textcolor{gregoriocolor}{A} & \textcolor{gregoriocolor}{B} & \textcolor{gregoriocolor}{C} & \textcolor{gregoriocolor}{D} & \textcolor{gregoriocolor}{E} & F & \textcolor{gregoriocolor}{F} & \textcolor{gregoriocolor}{G} & \textcolor{gregoriocolor}{H} & \textcolor{gregoriocolor}{M} & \textcolor{gregoriocolor}{N} & \textcolor{gregoriocolor}{P} \\
 23 & 24 & 25 & 26 & 27 & 28 & 28 & 29 & 30 & 1 & 2 & 3 \\
\end{tabularx}

\begin{paracol}{2}
\selectlanguage{latin}
\lettrine[lines=1]{O}{ctava} sancti Joánnis, Apóstoli et Evangelístæ.
\switchcolumn
\selectlanguage{english}
\lettrine[lines=1]{T}{he} Octave of St. John, apostle and evangelist.
\switchcolumn*
\selectlanguage{latin}
Romæ, via Appia, natális sancti Anthéri, Papæ 
 et Mártyris; qui sub Júlio Maximíno passus est, et in cœmetério Callísti 
 sepúltus.
\switchcolumn
\selectlanguage{english}
At Rome, on the Appian Way, the birthday of Pope St. Anterus, who suffered 
 under Julius Maximinus, and was buried in the cemetery of Callistus.
\switchcolumn*
\selectlanguage{latin}
Viénnæ, in Gállia, sancti Floréntii Epíscopi, 
 qui, témpore Galliéni Imperatóris, in exsílium relegátus, illic martyrium 
 consummávit.
\switchcolumn
\selectlanguage{english}
At Vienne in France, St. Florentius, bishop, who was sent into exile and who 
 was martyred in the time of Emperor Gallienus.
\switchcolumn*
\selectlanguage{latin}
Apud civitátem Aulánam, in Palæstína, pássio 
 sancti Petri, qui crucis supplício interémptus est.
\switchcolumn
\selectlanguage{english}
In the city of Aulona in Palestine, the crucifixion of St. Peter.
\switchcolumn*
\selectlanguage{latin}
In Hellespónto sanctórum Mártyrum Cyríni, Primi et Theógenis.
\switchcolumn
\selectlanguage{english}
In the Hellespont, the holy martyrs Cyrinus, Primus, and Theogenes.
\switchcolumn*
\selectlanguage{latin}
Cæsaréæ, in Cappadócia, sancti Górdii 
 Centuriónis, Mártyris; de cujus láudibus exstat præclára Basilíi Magni 
 orátio, in ejus die festo hábita.
\switchcolumn
\selectlanguage{english}
At Caesarea in Cappadocia, St. Gordius, centurion, in whose praise is extant 
 a celebrated discourse delivered by St. Basil the Great on the day of his 
 festival.
\switchcolumn*
\selectlanguage{latin}
In Cilícia sanctórum Mártyrum Zósimi, et 
 Athanásii Commentariénsis.
\switchcolumn
\selectlanguage{english}
In Cilicia, the holy martyrs Zosimus and the notary Athanasius.
\switchcolumn*
\selectlanguage{latin}
Item sanctórum Theopémpti et Theónæ, qui, in 
 persecutióne Diocletiáni, illústre martyrium obiérunt.
\switchcolumn
\selectlanguage{english}
Also, the Saints Theopemptus and Theonas, who suffered a glorious martyrdom 
 in the persecution of Diocletian.
\switchcolumn*
\selectlanguage{latin}
Patávii sancti Daniélis Mártyris.
\switchcolumn
\selectlanguage{english}
At Padua, St. Daniel, martyr.
\switchcolumn*
\selectlanguage{latin}
Lutétiæ Parisiórum sanctæ Genovéfæ Vírginis, 
 quæ, a beáto Germáno, Antisiodorénsi Epíscopo, Christo dicáta, admirándis 
 virtútibus et miráculis cláruit.
\switchcolumn
\selectlanguage{english}
At Paris, St. Genevieve, virgin, who was consecrated to Christ by St. 
 Germanus, bishop of Auxerre, and who became famous for her admirable virtues 
 and miracles.
\switchcolumn*
\selectlanguage{latin}
\end{paracol}


% ---- martyrology/mart01/mart0104.htm
\needspace{10\baselineskip}
\begin{paracol}{2}
\selectlanguage{latin}
\begin{center}{\color{gregoriocolor} Prídie Nonas Januárii. 
 Luna\dots\ }\end{center}
\switchcolumn
\selectlanguage{english}
\begin{center}{\color{gregoriocolor} The   4\textsuperscript{th} Day of 
 January. The\dots\ Day of the Moon.}\end{center}
\end{paracol}

\noindent\begin{tabularx}{\linewidth}{*{19}{>{\centering\arraybackslash}X}}
 \textcolor{gregoriocolor}{a} & \textcolor{gregoriocolor}{b} & \textcolor{gregoriocolor}{c} & \textcolor{gregoriocolor}{d} & \textcolor{gregoriocolor}{e} & \textcolor{gregoriocolor}{f} & \textcolor{gregoriocolor}{g} & \textcolor{gregoriocolor}{h} & \textcolor{gregoriocolor}{i} & \textcolor{gregoriocolor}{k} & \textcolor{gregoriocolor}{l} & \textcolor{gregoriocolor}{m} & \textcolor{gregoriocolor}{n} & \textcolor{gregoriocolor}{p} & \textcolor{gregoriocolor}{q} & \textcolor{gregoriocolor}{r} & \textcolor{gregoriocolor}{s} & \textcolor{gregoriocolor}{t} & \textcolor{gregoriocolor}{u} \\
 5 & 6 & 7 & 8 & 9 & 10 & 11 & 12 & 13 & 14 & 15 & 16 & 17 & 18 & 19 & 20 & 21 & 22 & 23 \\
\end{tabularx}
\vspace{0.5\baselineskip}
\noindent\begin{tabularx}{\linewidth}{*{12}{>{\centering\arraybackslash}X}}
 \textcolor{gregoriocolor}{A} & \textcolor{gregoriocolor}{B} & \textcolor{gregoriocolor}{C} & \textcolor{gregoriocolor}{D} & \textcolor{gregoriocolor}{E} & F & \textcolor{gregoriocolor}{F} & \textcolor{gregoriocolor}{G} & \textcolor{gregoriocolor}{H} & \textcolor{gregoriocolor}{M} & \textcolor{gregoriocolor}{N} & \textcolor{gregoriocolor}{P} \\
 24 & 25 & 26 & 27 & 28 & 29 & 29 & 30 & 1 & 2 & 3 & 4 \\
\end{tabularx}

\begin{paracol}{2}
\selectlanguage{latin}
\lettrine[lines=1]{O}{ctáva} sanctórum Innocéntium Mártyrum.
\switchcolumn
\selectlanguage{english}
\lettrine[lines=1]{T}{he} Octave of the Holy Innocents.
\switchcolumn*
\selectlanguage{latin}
In Creta natális sancti Titi, qui, ab Apóstolo Paulo Epíscopus Creténsium 
 ordinátus, et, post prædicatiónis offícium 
 fidelíssime consummátum, finem beátum adéptus, in ea sepúltus est Ecclésia, 
 ubi a beáto Apóstolo dignus miníster fúerat constitútus. Ipsíus tamen 
 festívitas octávo Idus Februárii celebrátur.
\switchcolumn
\selectlanguage{english}
In Crete, the birthday of St. Titus, who was consecrated bishop of that 
 island by the apostle St. Paul. After having faithfully performed the 
 duty of preaching the Gospel, he reached the end of his saintly life, and 
 was buried in the church of which he had been made a worthy minister by the 
 holy apostle.
\switchcolumn*
\selectlanguage{latin}
Romæ sanctórum Mártyrum Prisci Presbyteri, et 
 Priscilliáni Clérici, ac Benedíctæ, religiósæ féminæ; qui, témpore 
 impiíssimi Juliáni, gládio martyrium complevérunt.
\switchcolumn
\selectlanguage{english}
At Rome, in the reign of the impious Julian, the holy martyrs Priscus, a 
 priest, Priscillian, a cleric; and Benedicta, a religious woman, whose 
 martyrdom was ended by the sword.
\switchcolumn*
\selectlanguage{latin}
Item Romæ beátæ Dafrósæ, uxóris sancti Flaviáni 
 Mártyris, ac matris sanctárum Bibiánæ et Demétriæ, Vírginum et Mártyrum; quæ, 
 post interfectiónem viri sui, primum exsílio relegáta, deínde, sub præfáto 
 Príncipe, cápite plexa est.
\switchcolumn
\selectlanguage{english}
Also at Rome, under the same emperor, blessed Dafrosa, wife of the martyr 
 St. Flavian, and mother of Saints Bibiana and Demetria, virgin martyrs. 
 After her husband had been killed, she was first banished and then beheaded.
\switchcolumn*
\selectlanguage{latin}
Bonóniæ sanctórum Hermétis, Aggǽi et Caji 
 Mártyrum, qui sub Maximiáno Imperatóre passi sunt.
\switchcolumn
\selectlanguage{english}
At Bologna, the Saints Hermes, Aggaeus, and Caius, martyrs, who suffered 
 under Emperor Maximian.
\switchcolumn*
\selectlanguage{latin}
Adruméti, in Africa, commemorátio sancti Mávili 
 Mártyris, qui in persecutióne Sevéri Imperatóris, a sævíssimo Præside 
 Scápula damnátus ad béstias, martyrii corónam accépit.
\switchcolumn
\selectlanguage{english}
At Adrumetum in Africa, in the persecution of Severus, the commemoration of 
 St. Mavilus, martyr, who, being condemned by the very cruel governor Scapula 
 to be devoured by wild beasts, received the crown of martyrdom.
\switchcolumn*
\selectlanguage{latin}
Item in Africa præclarissimórum Mártyrum 
 Aquilíni, Gémini, Eugénii, Marciáni, Quincti, Theódoti et Tryphónis.
\switchcolumn
\selectlanguage{english}
Also in Africa, the most renowned martyrs Aquilinus, Geminus, Eugenius, 
 Marcian, Quinctus, Theodotus, and Tryphon.
\switchcolumn*
\selectlanguage{latin}
Apud Língonas, in Gállia, sancti Gregórii 
 Epíscopi, miráculis clari.
\switchcolumn
\selectlanguage{english}
At Langres in France, St. Gregory, a bishop renowned for miracles.
\switchcolumn*
\selectlanguage{latin}
Rhemis, in Gállia, sancti Rigobérti, Epíscopi et Confessóris.
\switchcolumn
\selectlanguage{english}
At Rheims in France, St. Rigobertus, bishop and confessor.
\switchcolumn*
\selectlanguage{latin}
\end{paracol}


% ---- martyrology/mart01/mart0105.htm
\needspace{10\baselineskip}
\begin{paracol}{2}
\selectlanguage{latin}
\begin{center}{\color{gregoriocolor} Nonis Januárii. 
 Luna\dots\ }\end{center}
\switchcolumn
\selectlanguage{english}
\begin{center}{\color{gregoriocolor} The   5\textsuperscript{th} Day of 
 January. The\dots\ Day of the Moon.}\end{center}
\end{paracol}

\noindent\begin{tabularx}{\linewidth}{*{19}{>{\centering\arraybackslash}X}}
 \textcolor{gregoriocolor}{a} & \textcolor{gregoriocolor}{b} & \textcolor{gregoriocolor}{c} & \textcolor{gregoriocolor}{d} & \textcolor{gregoriocolor}{e} & \textcolor{gregoriocolor}{f} & \textcolor{gregoriocolor}{g} & \textcolor{gregoriocolor}{h} & \textcolor{gregoriocolor}{i} & \textcolor{gregoriocolor}{k} & \textcolor{gregoriocolor}{l} & \textcolor{gregoriocolor}{m} & \textcolor{gregoriocolor}{n} & \textcolor{gregoriocolor}{p} & \textcolor{gregoriocolor}{q} & \textcolor{gregoriocolor}{r} & \textcolor{gregoriocolor}{s} & \textcolor{gregoriocolor}{t} & \textcolor{gregoriocolor}{u} \\
 6 & 7 & 8 & 9 & 10 & 11 & 12 & 13 & 14 & 15 & 16 & 17 & 18 & 19 & 20 & 21 & 22 & 23 & 24 \\
\end{tabularx}
\vspace{0.5\baselineskip}
\noindent\begin{tabularx}{\linewidth}{*{12}{>{\centering\arraybackslash}X}}
 \textcolor{gregoriocolor}{A} & \textcolor{gregoriocolor}{B} & \textcolor{gregoriocolor}{C} & \textcolor{gregoriocolor}{D} & \textcolor{gregoriocolor}{E} & F & \textcolor{gregoriocolor}{F} & \textcolor{gregoriocolor}{G} & \textcolor{gregoriocolor}{H} & \textcolor{gregoriocolor}{M} & \textcolor{gregoriocolor}{N} & \textcolor{gregoriocolor}{P} \\
 25 & 26 & 27 & 28 & 29 & 30 & 30 & 1 & 2 & 3 & 4 & 5 \\
\end{tabularx}

\begin{paracol}{2}
\selectlanguage{latin}
\lettrine[lines=1]{V}{igília} Epiphaníæ Dómini.
\switchcolumn
\selectlanguage{english}
\lettrine[lines=1]{T}{he} Vigil of the Epiphany of our Lord.
\switchcolumn*
\selectlanguage{latin}
Romæ sancti Telésphori, Papæ et Mártyris; qui, 
 sub Antoníno Pio, post multos labóres, pro Christi confessióne, illústre 
 martyrium duxit.
\switchcolumn
\selectlanguage{english}
At Rome, in the time of Antoninus Pius, St. Telesphorus, pope, who, after 
 many sufferings for the confession of Christ, underwent a glorious 
 martyrdom.
\switchcolumn*
\selectlanguage{latin}
In Anglia natális sancti Eduárdi, Regis Anglórum et Confessóris; qui virtúte 
 castitátis et grátia miraculórum fuit insígnis. Ejus autem festívitas, 
 ex decréto Innocéntii Papæ Undécimi, tértio 
 Idus Octóbris, quo die sacrum ejus corpus translátum fuit, potíssimum 
 celebrátur.
\switchcolumn
\selectlanguage{english}
In England, St. Edward, king and confessor, illustrious by the virtue of 
 chastity and the gift of miracles. His feast, by order of Pope 
 Innocent XI, is celebrated on the 13th of October, on which day his holy 
 body was transferred.
\switchcolumn*
\selectlanguage{latin}
In Ægypto commemorátio plurimórum sanctórum 
 Mártyrum, qui in Thebáide, sub persecutióne Diocletiáni, divérso tormentórum 
 génere cæsi sunt.
\switchcolumn
\selectlanguage{english}
In Egypt, during the persecution of Diocletian, the commemoration of many 
 holy martyrs who were put to death in Thebais by various kinds of torments.
\switchcolumn*
\selectlanguage{latin}
Antiochíæ sancti Simeónis Mónachi, qui, multos 
 annos in colúmna stans vixit, unde et Stylítæ cognómen accépit; cujus vita 
 et conversátio éxstitit admirábilis.
\switchcolumn
\selectlanguage{english}
At Antioch, St. Simeon, monk, admirable both for his life and for his 
 conversation. He lived for many years standing on a pillar, and was 
 for that reason called Stylites.
\switchcolumn*
\selectlanguage{latin}
Romæ sanctæ Æmiliánæ Vírginis, ámitæ sancti 
 Gregórii Papæ; quæ, vocánte Tharsílla soróre, quæ ad Deum præcésserat, hac 
 ipsa die migrávit ad Dóminum.
\switchcolumn
\selectlanguage{english}
At Rome, the holy virgin Emiliana, aunt of Pope St. Gregory. Being 
 called to God by her sister Tharsilla, who had preceded her, she departed to 
 heaven on this day.
\switchcolumn*
\selectlanguage{latin}
Alexandríæ sanctæ Syncléticæ Vírginis, cujus 
 res præcláre gestas sanctus Athanásius monuméntis litterárum commendávit.
\switchcolumn
\selectlanguage{english}
At Alexandria, St. Syncletica, whose noble deeds have been recorded by St. 
 Athanasius.
\switchcolumn*
\selectlanguage{latin}
In Ægypto sanctæ Apollináris Vírginis.
\switchcolumn
\selectlanguage{english}
In Egypt, St. Apollinaris, virgin.
\switchcolumn*
\selectlanguage{latin}
\end{paracol}


% ---- martyrology/mart01/mart0106.htm
\needspace{10\baselineskip}
\begin{paracol}{2}
\selectlanguage{latin}
\begin{center}{\color{gregoriocolor} Octávo Idus Januárii. 
 Luna\dots\ }\end{center}
\switchcolumn
\selectlanguage{english}
\begin{center}{\color{gregoriocolor} The 6\textsuperscript{th} Day of 
 January. The\dots\ Day of the Moon.}\end{center}
\end{paracol}

\noindent\begin{tabularx}{\linewidth}{*{19}{>{\centering\arraybackslash}X}}
 \textcolor{gregoriocolor}{a} & \textcolor{gregoriocolor}{b} & \textcolor{gregoriocolor}{c} & \textcolor{gregoriocolor}{d} & \textcolor{gregoriocolor}{e} & \textcolor{gregoriocolor}{f} & \textcolor{gregoriocolor}{g} & \textcolor{gregoriocolor}{h} & \textcolor{gregoriocolor}{i} & \textcolor{gregoriocolor}{k} & \textcolor{gregoriocolor}{l} & \textcolor{gregoriocolor}{m} & \textcolor{gregoriocolor}{n} & \textcolor{gregoriocolor}{p} & \textcolor{gregoriocolor}{q} & \textcolor{gregoriocolor}{r} & \textcolor{gregoriocolor}{s} & \textcolor{gregoriocolor}{t} & \textcolor{gregoriocolor}{u} \\
 7 & 8 & 9 & 10 & 11 & 12 & 13 & 14 & 15 & 16 & 17 & 18 & 19 & 20 & 21 & 22 & 23 & 24 & 25 \\
\end{tabularx}
\vspace{0.5\baselineskip}
\noindent\begin{tabularx}{\linewidth}{*{12}{>{\centering\arraybackslash}X}}
 \textcolor{gregoriocolor}{A} & \textcolor{gregoriocolor}{B} & \textcolor{gregoriocolor}{C} & \textcolor{gregoriocolor}{D} & \textcolor{gregoriocolor}{E} & F & \textcolor{gregoriocolor}{F} & \textcolor{gregoriocolor}{G} & \textcolor{gregoriocolor}{H} & \textcolor{gregoriocolor}{M} & \textcolor{gregoriocolor}{N} & \textcolor{gregoriocolor}{P} \\
 26 & 27 & 28 & 29 & 30 & 1 & 1 & 2 & 3 & 4 & 5 & 6 \\
\end{tabularx}

\begin{paracol}{2}
\selectlanguage{latin}
\lettrine[lines=1]{E}{piphanía} Domini.
\switchcolumn
\selectlanguage{english}
\lettrine[lines=1]{T}{he} Epiphany of our Lord.
\switchcolumn*
\selectlanguage{latin}
Floréntiæ natális sancti Andréæ Corsíni, civis 
 Florentíni, ex Ordine Carmelitárum, Epíscopi Fæsuláni et Confessóris; quem, 
 miráculis clarum, Urbánus Papa Octávus in Sanctórum númerum rétulit. Ejus autem festívitas recólitur prídie nonas Februárii.
\switchcolumn
\selectlanguage{english}
At Florence, St. Andrew Corsini, a Florentine Carmelite and bishop of 
 Fiesole. Being celebrated for miracles, he was ranked among the saints 
 by Urban VIII. His festival is kept on the 4th of February.
\switchcolumn*
\selectlanguage{latin}
Barcinóne, in Hispánia, item natális sancti Raymúndi de Pénafort, ex Ordine 
 Prædicatórum, Confessóris, doctrína et 
 sanctitáte célebris. Ipsíus vero festum décimo Kaléndas Februárii 
 celebrátur.
\switchcolumn
\selectlanguage{english}
At Barcelona in Spain, St. Raymond of Pennafort, of the Order of Preachers, 
 celebrated for sanctity and learning. His festival is kept on the 23rd 
 of this month.
\switchcolumn*
\selectlanguage{latin}
In Africa commemorátio plurimórum sanctórum Mártyrum, qui, in persecutióne 
 Sevéri, ad palum ligáti sunt et igne consúmpti.
\switchcolumn
\selectlanguage{english}
In Africa, the commemoration of many holy martyrs who were burned at the 
 stake in the persecution of Severus.
\switchcolumn*
\selectlanguage{latin}
In território Rheménsi pássio sanctæ Macræ 
 Vírginis, quæ, in persecutióne Diocletiáni, jubénte Rictiováro Præside, cum 
 in ignem esset præcipitáta et permansísset illæsa, dehinc, mamíllis 
 abscíssis et squalóre cárceris afflícta, super testas étiam acutíssimas ac 
 prunas volutáta, tandem orans migrávit ad Dóminum.
\switchcolumn
\selectlanguage{english}
In the diocese of Rheims, the martyrdom of St. Macra, virgin, who, in the 
 persecution of Diocletian, was cast into the fire by order of the governor 
 Rictiovarus. As she remained uninjured, she had her breasts cut away, 
 was imprisoned in a foul dungeon, rolled upon broken earthenware and burning 
 coals, and finally she gave up her soul while engaged in prayer.
\switchcolumn*
\selectlanguage{latin}
Rhédonis, in Gállia, sancti Melánii, Epíscopi 
 et Confessóris; qui, post innumerabílium signa virtútum, júgiter cælo 
 inténtus, gloriósus migrávit a sæculo.
\switchcolumn
\selectlanguage{english}
At Rennes in France, St. Melanius, bishop and confessor, who, after a life 
 remarkable for innumerable virtues, with his thoughts constantly fixed on 
 heaven, gloriously departed from this world.
\switchcolumn*
\selectlanguage{latin}
Geris, in Ægypto, sancti Nilammónis reclúsi, 
 qui, dum ad Episcopátum traherétur invítus, in oratióne spíritum Deo 
 réddidit.
\switchcolumn
\selectlanguage{english}
At Geris in Egypt, St. Nilammon, anchoret, who, while he was carried to a 
 bishopric against his will, gave up his soul to God in prayer.
\switchcolumn*
\selectlanguage{latin}
\end{paracol}


% ---- martyrology/mart01/mart0107.htm
\needspace{10\baselineskip}
\begin{paracol}{2}
\selectlanguage{latin}
\begin{center}{\color{gregoriocolor} Séptimo Idus Januárii. 
 Luna\dots\ }\end{center}
\switchcolumn
\selectlanguage{english}
\begin{center}{\color{gregoriocolor} The 7\textsuperscript{th} Day of 
 January. The\dots\ Day of the Moon.}\end{center}
\end{paracol}

\noindent\begin{tabularx}{\linewidth}{*{19}{>{\centering\arraybackslash}X}}
 \textcolor{gregoriocolor}{a} & \textcolor{gregoriocolor}{b} & \textcolor{gregoriocolor}{c} & \textcolor{gregoriocolor}{d} & \textcolor{gregoriocolor}{e} & \textcolor{gregoriocolor}{f} & \textcolor{gregoriocolor}{g} & \textcolor{gregoriocolor}{h} & \textcolor{gregoriocolor}{i} & \textcolor{gregoriocolor}{k} & \textcolor{gregoriocolor}{l} & \textcolor{gregoriocolor}{m} & \textcolor{gregoriocolor}{n} & \textcolor{gregoriocolor}{p} & \textcolor{gregoriocolor}{q} & \textcolor{gregoriocolor}{r} & \textcolor{gregoriocolor}{s} & \textcolor{gregoriocolor}{t} & \textcolor{gregoriocolor}{u} \\
 8 & 9 & 10 & 11 & 12 & 13 & 14 & 15 & 16 & 17 & 18 & 19 & 20 & 21 & 22 & 23 & 24 & 25 & 26 \\
\end{tabularx}
\vspace{0.5\baselineskip}
\noindent\begin{tabularx}{\linewidth}{*{12}{>{\centering\arraybackslash}X}}
 \textcolor{gregoriocolor}{A} & \textcolor{gregoriocolor}{B} & \textcolor{gregoriocolor}{C} & \textcolor{gregoriocolor}{D} & \textcolor{gregoriocolor}{E} & F & \textcolor{gregoriocolor}{F} & \textcolor{gregoriocolor}{G} & \textcolor{gregoriocolor}{H} & \textcolor{gregoriocolor}{M} & \textcolor{gregoriocolor}{N} & \textcolor{gregoriocolor}{P} \\
 27 & 28 & 29 & 30 & 1 & 2 & 2 & 3 & 4 & 5 & 6 & 7 \\
\end{tabularx}

\begin{paracol}{2}
\selectlanguage{latin}
\lettrine[lines=1]{R}{elátio} púeri Jesu de Ægypto.
\switchcolumn
\selectlanguage{english}
\lettrine[lines=1]{T}{he} return of the Child Jesus from Egypt.
\switchcolumn*
\selectlanguage{latin}
Nicomedíæ natális beáti Luciáni, Ecclésiæ 
 Antiochénæ Presbyteri et Mártyris; qui, satis clarus doctrína et eloquéntia, 
 passus est, ob Christi confessiónem, in persecutióne Galérii Maximiáni, 
 sepultúsque est Helenópoli, in Bithynia. Ipsíus autem laudes sanctus 
 Joánnes Chrysóstomus celebrávit.
\switchcolumn
\selectlanguage{english}
The birthday of blessed Lucian, a priest of the Church of Antioch and 
 martyr, who was distinguished for his learning and eloquence. He 
 suffered at Nicomedia for the confession of Christ, in the persecution of 
 Galerius Maximian, and was buried at Helenopolis, in Bithynia. His 
 praises have been proclaimed by St. John Chrysostom.
\switchcolumn*
\selectlanguage{latin}
Antiochíæ sancti Cleri Diáconi, qui pro 
 confessiónis glória, sépties tortus ac diu macerátus in cárcere, ad últimum, 
 gládio decollátus, martyrium consummávit.
\switchcolumn
\selectlanguage{english}
At Antioch, St. Clerus, deacon, who, for having professed faith in Christ, 
 was seven times tortured, kept in prison a long while, and at length his 
 martyrdom was ended by decapitation.
\switchcolumn*
\selectlanguage{latin}
In civitáte Heracléa sanctórum Mártyrum Felícis et Januárii.
\switchcolumn
\selectlanguage{english}
In the city of Heraclea, the holy martyrs Felix and Januarius.
\switchcolumn*
\selectlanguage{latin}
Eódem die sancti Juliáni Mártyris.
\switchcolumn
\selectlanguage{english}
The same day, St. Julian, martyr.
\switchcolumn*
\selectlanguage{latin}
In Dánia sancti Canúti, Regis et Mártyris.
\switchcolumn
\selectlanguage{english}
In Denmark, St. Canute, king and martyr.
\switchcolumn*
\selectlanguage{latin}
Papíæ sancti Crispíni, Epíscopi et Confessóris.
\switchcolumn
\selectlanguage{english}
At Pavia, St. Crispin, bishop and confessor.
\switchcolumn*
\selectlanguage{latin}
In Dácia sancti Nicétæ Epíscopi, qui feras et 
 bárbaras gentes, Evangélii prædicatióne, mites réddidit ac mansuétas.
\switchcolumn
\selectlanguage{english}
In Dacia, St. Nicetas, bishop, who made fierce and barbarous nations humane 
 and meek by preaching the Gospel to them.
\switchcolumn*
\selectlanguage{latin}
In Ægypto beáti Theodóri Mónachi, qui, témpore 
 Constantíni Magni, flóruit sanctitáte; cujus méminit sanctus Athanásius in 
 vita sancti Antónii.
\switchcolumn
\selectlanguage{english}
In Egypt, St. Theodore, a saintly monk, who flourished in the time of 
 Constantine the Great. He is mentioned by St. Athanasius in his Life 
 of St. Anthony.
\switchcolumn*
\selectlanguage{latin}
\end{paracol}


% ---- martyrology/mart01/mart0108.htm
\needspace{10\baselineskip}
\begin{paracol}{2}
\selectlanguage{latin}
\begin{center}{\color{gregoriocolor} Sexto Idus Januárii. 
 Luna\dots\ }\end{center}
\switchcolumn
\selectlanguage{english}
\begin{center}{\color{gregoriocolor} The 8\textsuperscript{th} Day of 
 January. The\dots\ Day of the Moon.}\end{center}
\end{paracol}

\noindent\begin{tabularx}{\linewidth}{*{19}{>{\centering\arraybackslash}X}}
 \textcolor{gregoriocolor}{a} & \textcolor{gregoriocolor}{b} & \textcolor{gregoriocolor}{c} & \textcolor{gregoriocolor}{d} & \textcolor{gregoriocolor}{e} & \textcolor{gregoriocolor}{f} & \textcolor{gregoriocolor}{g} & \textcolor{gregoriocolor}{h} & \textcolor{gregoriocolor}{i} & \textcolor{gregoriocolor}{k} & \textcolor{gregoriocolor}{l} & \textcolor{gregoriocolor}{m} & \textcolor{gregoriocolor}{n} & \textcolor{gregoriocolor}{p} & \textcolor{gregoriocolor}{q} & \textcolor{gregoriocolor}{r} & \textcolor{gregoriocolor}{s} & \textcolor{gregoriocolor}{t} & \textcolor{gregoriocolor}{u} \\
 9 & 10 & 11 & 12 & 13 & 14 & 15 & 16 & 17 & 18 & 19 & 20 & 21 & 22 & 23 & 24 & 25 & 26 & 27 \\
\end{tabularx}
\vspace{0.5\baselineskip}
\noindent\begin{tabularx}{\linewidth}{*{12}{>{\centering\arraybackslash}X}}
 \textcolor{gregoriocolor}{A} & \textcolor{gregoriocolor}{B} & \textcolor{gregoriocolor}{C} & \textcolor{gregoriocolor}{D} & \textcolor{gregoriocolor}{E} & F & \textcolor{gregoriocolor}{F} & \textcolor{gregoriocolor}{G} & \textcolor{gregoriocolor}{H} & \textcolor{gregoriocolor}{M} & \textcolor{gregoriocolor}{N} & \textcolor{gregoriocolor}{P} \\
 28 & 29 & 30 & 1 & 2 & 3 & 3 & 4 & 5 & 6 & 7 & 8 \\
\end{tabularx}

\begin{paracol}{2}
\selectlanguage{latin}
\lettrine[lines=2]{V}{enétiis} deposítio sancti Lauréntii Justiniáni, primi Patriárchæ 
 urbis ejúsdem et Confessóris; quem, doctrína et supérnis divínæ sapiéntiæ 
 charismátibus copiosíssime replétum, Alexánder Octávus, Póntifex Máximus, in 
 Sanctórum númerum rétulit. Ipsíus autem festívitas Nonis Septémbris, 
 quo die Cáthedram pontificálem ascéndit, potíssimum celebrátur.
\switchcolumn
\selectlanguage{english}
\lettrine[lines=2]{A}{t} Venice, the death of St. Lawrence Justinian, confessor, first patriarch 
 of that city. Eminent for learning, and abundantly filled with the 
 heavenly gifts of divine wisdom, he was ranked among the saints by Alexander 
 VIII. He is again mentioned on the 5th of September, on which day he 
 ascended the pontifical throne.
\switchcolumn*
\selectlanguage{latin}
Bellóvaci, in Gálliis, sanctórum Mártyrum 
 Luciáni Presbyteri, Maximiáni et Juliáni. Horum duo últimi a persecutóribus 
 gládio perémpti sunt; beátum autem Luciánus, qui, una cum sancto Dionysio, 
 in Gálliam vénerat, et ipse, post nímiam cædem, cum Christi nomen viva voce 
 confitéri non timuísset, priórum senténtiam excépit.
\switchcolumn
\selectlanguage{english}
At Beauvais in France, the holy martyrs Lucian, priest, Maximian and Julian. 
 The last two were killed with the sword by the persecutors; but blessed 
 Lucian, who had come to France with St. Denis, after the slaughter of his 
 companions, not fearing to confess the Name of Christ openly, received the 
 same sentence of death.
\switchcolumn*
\selectlanguage{latin}
In Libya sanctórum Mártyrum Theóphili Diáconi, et Helládii, qui, primo 
 laniáti ac téstulis peracútis compúncti, demum, in ignem conjécti, ánimas 
 Deo reddidérunt.
\switchcolumn
\selectlanguage{english}
In Libya, the holy martyrs Theophilus, deacon, and Helladius, who, after 
 having their bodies lacerated and cut with sharp pieces of earthenware, were 
 cast into the fire, and rendered their souls unto God.
\switchcolumn*
\selectlanguage{latin}
Augustodúni sancti Eugeniáni Mártyris.
\switchcolumn
\selectlanguage{english}
At Autun, St. Eugenian, martyr.
\switchcolumn*
\selectlanguage{latin}
Hierápoli, in Asia, sancti Apollináris Epíscopi, qui, sub Marco Antoníno 
 Vero, sanctitáte atque doctrína flóruit.
\switchcolumn
\selectlanguage{english}
At Hierapolis in Asia, St. Apollinaris, bishop, who was conspicuous for 
 sanctity and learning under Marucs Antoninus Verus.
\switchcolumn*
\selectlanguage{latin}
Neápoli, in Campánia, natális sancti Severíni Epíscopi, qui fuit frater 
 beáti Victoríni Mártyris; et, post multárum virtútum perpetratiónem, plenus 
 sanctitáte quiévit.
\switchcolumn
\selectlanguage{english}
At Naples in Campania, the birthday of the bishop St. Severin, brother to 
 the blessed martyr Victorinus, who, after working many miracles, died, 
 replenished with virtues and merits.
\switchcolumn*
\selectlanguage{latin}
Metis, in Gállia, sancti Patiéntis Epíscopi.
\switchcolumn
\selectlanguage{english}
At Metz in France, St. Patiens, bishop.
\switchcolumn*
\selectlanguage{latin}
Papíæ sancti Máximi, Epíscopi et Confessóris.
\switchcolumn
\selectlanguage{english}
At Pavia, St. Maximus, bishop and confessor.
\switchcolumn*
\selectlanguage{latin}
Ratisbónæ, in Bavária, sancti Erhárdi Epíscopi.
\switchcolumn
\selectlanguage{english}
At Ratisbon in Bavaria, St. Erhard, bishop.
\switchcolumn*
\selectlanguage{latin}
Apud Nóricos sancti Severíni Abbátis, qui apud 
 eam gentem Evangélium propagávit, et Noricórum dictus est Apóstolus. 
 Ejus corpus ad Lucullánum prope Neápolim, in Campánia, divínitus delátum, 
 inde póstea ad monastérium sancti Severíni translátum est.
\switchcolumn
\selectlanguage{english}
Among the inhabitants of Noricum (now Austria), the abbot St. Severin, who 
 propagated the Gospel in that country, and is called its apostle. By 
 divine power his body was carried to Lucullano, near Naples, and thence 
 transferred to the monastery of St. Severin.
\switchcolumn*
\selectlanguage{latin}
\end{paracol}


% ---- martyrology/mart01/mart0109.htm
\needspace{10\baselineskip}
\begin{paracol}{2}
\selectlanguage{latin}
\begin{center}{\color{gregoriocolor} Quinto Idus Januárii. 
 Luna\dots\ }\end{center}
\switchcolumn
\selectlanguage{english}
\begin{center}{\color{gregoriocolor} The 9\textsuperscript{th} Day of 
 January. The\dots\ Day of the Moon.}\end{center}
\end{paracol}

\noindent\begin{tabularx}{\linewidth}{*{19}{>{\centering\arraybackslash}X}}
 \textcolor{gregoriocolor}{a} & \textcolor{gregoriocolor}{b} & \textcolor{gregoriocolor}{c} & \textcolor{gregoriocolor}{d} & \textcolor{gregoriocolor}{e} & \textcolor{gregoriocolor}{f} & \textcolor{gregoriocolor}{g} & \textcolor{gregoriocolor}{h} & \textcolor{gregoriocolor}{i} & \textcolor{gregoriocolor}{k} & \textcolor{gregoriocolor}{l} & \textcolor{gregoriocolor}{m} & \textcolor{gregoriocolor}{n} & \textcolor{gregoriocolor}{p} & \textcolor{gregoriocolor}{q} & \textcolor{gregoriocolor}{r} & \textcolor{gregoriocolor}{s} & \textcolor{gregoriocolor}{t} & \textcolor{gregoriocolor}{u} \\
 10 & 11 & 12 & 13 & 14 & 15 & 16 & 17 & 18 & 19 & 20 & 21 & 22 & 23 & 24 & 25 & 26 & 27 & 28 \\
\end{tabularx}
\vspace{0.5\baselineskip}
\noindent\begin{tabularx}{\linewidth}{*{12}{>{\centering\arraybackslash}X}}
 \textcolor{gregoriocolor}{A} & \textcolor{gregoriocolor}{B} & \textcolor{gregoriocolor}{C} & \textcolor{gregoriocolor}{D} & \textcolor{gregoriocolor}{E} & F & \textcolor{gregoriocolor}{F} & \textcolor{gregoriocolor}{G} & \textcolor{gregoriocolor}{H} & \textcolor{gregoriocolor}{M} & \textcolor{gregoriocolor}{N} & \textcolor{gregoriocolor}{P} \\
 29 & 30 & 1 & 2 & 3 & 4 & 4 & 5 & 6 & 7 & 8 & 9 \\
\end{tabularx}

\begin{paracol}{2}
\selectlanguage{latin}
\lettrine[lines=2]{A}{ntiochíæ,} sub Diocletiáno et Maximiáno, 
 natális sanctórum Juliáni Mártyris, et Basilíssæ Vírginis, ipsíus Juliáni 
 uxóris. Hæc, virginitáte cum viro suo serváta, in pace vitam finívit; 
 Juliánus vero (postquam multitúdo Sacerdótum et Ministrórum Ecclésiæ 
 Christi, quæ, propter immanitátem persecutiónis, ad eos confúgerat, igne 
 cremáta est), Marciáni Præsidis jussu, plúrimis torméntis cruciátus, 
 capitálem senténtiam accépit. Cum ipso étiam Antónius Présbyter, et Anastásius, quem idem Juliánus, a morte suscitátum, grátiæ Christi partícipem fécerat, et Celsus puer cum hujus matre Marcionílla, ac septem 
 fratres, aliíque plúrimi passi sunt.
\switchcolumn
\selectlanguage{english}
\lettrine[lines=2]{A}{t} Antioch, in the reign of Diocletian and Maximian, the birthday of the 
 Saints Julian, martyr, and Basilissa, his virgin wife. She, having 
 lived in a state of virginity with her husband, reached the end of her days 
 in peace. But Julian, after the death by fire of a multitude of 
 priests and ministers of the Church of Christ, who had taken refuge in his 
 house from the severity of the persecution, was ordered by the governor 
 Marcian to be tormented in many ways and executed. With him there 
 suffered Anthony, a priest, and Anastasius, whom Julian raised from the 
 dead, and made partaker of the grace of Christ; also Celsus, a boy, with his 
 mother Marcionilla, seven brothers, and many others.
\switchcolumn*
\selectlanguage{latin}
Smyrnæ sanctórum Mártyrum Vitális, Revocáti et 
 Fortunáti.
\switchcolumn
\selectlanguage{english}
At Smyrna, the holy martyrs Vitalis, Revocatus, and Fortunatus.
\switchcolumn*
\selectlanguage{latin}
In Africa sanctórum Mártyrum Epictéti, Jucúndi, 
 Secúndi, Vitális, Felícis et aliórum septem.
\switchcolumn
\selectlanguage{english}
In Africa, the holy martyrs Epictetus, Jucundus, Secundus, Vitalis, Felix, 
 and seven others.
\switchcolumn*
\selectlanguage{latin}
In Mauritánia Cæsariénsi sanctæ Marciánæ 
 Vírginis, quæ, béstiis trádita, martyrium consummávit.
\switchcolumn
\selectlanguage{english}
In Algeria, St. Marciana, virgin, who received her martyrdom after being 
 condemned to the beasts.
\switchcolumn*
\selectlanguage{latin}
Sebáste, in Arménia, sancti Petri Epíscopi, fílii sanctórum Basilíi 
 et Emméliæ, atque fratris item sanctórum Basilíi Magni et Gregórii Nysséni 
 Episcopórum, ac Macrínæ Vírginis.
\switchcolumn
\selectlanguage{english}
At Sebaste in Armenia, St. Peter, bishop, the son of Saints Basil and 
 Emmelia, and also the brother of Saints Basil the Great, Gregory of Nyssa, 
 bishops, and Macrina, virgin.
\switchcolumn*
\selectlanguage{latin}
Ancónæ sancti Marcellíni Epíscopi, qui urbem 
 illam (ut sanctus Gregórius Papa scribit) divína virtúte ab incéndio 
 liberávit.
\switchcolumn
\selectlanguage{english}
At Ancona, St. Marcellinus, bishop, who, according to St. Gregory, 
 miraculously delivered that city from destruction by fire.
\switchcolumn*
\selectlanguage{latin}
\end{paracol}


% ---- martyrology/mart01/mart0110.htm
\needspace{10\baselineskip}
\begin{paracol}{2}
\selectlanguage{latin}
\begin{center}{\color{gregoriocolor} Quarto Idus Januárii. 
 Luna\dots\ }\end{center}
\switchcolumn
\selectlanguage{english}
\begin{center}{\color{gregoriocolor} The 10\textsuperscript{th} Day of 
 January. The\dots\ Day of the Moon.}\end{center}
\end{paracol}

\noindent\begin{tabularx}{\linewidth}{*{19}{>{\centering\arraybackslash}X}}
 \textcolor{gregoriocolor}{a} & \textcolor{gregoriocolor}{b} & \textcolor{gregoriocolor}{c} & \textcolor{gregoriocolor}{d} & \textcolor{gregoriocolor}{e} & \textcolor{gregoriocolor}{f} & \textcolor{gregoriocolor}{g} & \textcolor{gregoriocolor}{h} & \textcolor{gregoriocolor}{i} & \textcolor{gregoriocolor}{k} & \textcolor{gregoriocolor}{l} & \textcolor{gregoriocolor}{m} & \textcolor{gregoriocolor}{n} & \textcolor{gregoriocolor}{p} & \textcolor{gregoriocolor}{q} & \textcolor{gregoriocolor}{r} & \textcolor{gregoriocolor}{s} & \textcolor{gregoriocolor}{t} & \textcolor{gregoriocolor}{u} \\
 11 & 12 & 13 & 14 & 15 & 16 & 17 & 18 & 19 & 20 & 21 & 22 & 23 & 24 & 25 & 26 & 27 & 28 & 29 \\
\end{tabularx}
\vspace{0.5\baselineskip}
\noindent\begin{tabularx}{\linewidth}{*{12}{>{\centering\arraybackslash}X}}
 \textcolor{gregoriocolor}{A} & \textcolor{gregoriocolor}{B} & \textcolor{gregoriocolor}{C} & \textcolor{gregoriocolor}{D} & \textcolor{gregoriocolor}{E} & F & \textcolor{gregoriocolor}{F} & \textcolor{gregoriocolor}{G} & \textcolor{gregoriocolor}{H} & \textcolor{gregoriocolor}{M} & \textcolor{gregoriocolor}{N} & \textcolor{gregoriocolor}{P} \\
 30 & 1 & 2 & 3 & 4 & 5 & 5 & 6 & 7 & 8 & 9 & 10 \\
\end{tabularx}

\begin{paracol}{2}
\selectlanguage{latin}
\lettrine[lines=2]{I}{n} Thebáide natális beáti Pauli, primi Eremítæ, 
 Confessóris, qui, a sextodécimo ætátis suæ anno usque ad centésimum décimum 
 tértium, solus in erémo permánsit; cujus ánimam, inter Apostolórum et 
 Prophetárum choros, ad cælum ferri ab Angelis sanctus Antónius vidit. 
 Ejus autem festívitas décimo octávo Kaléndas Februárii celebrátur.
\switchcolumn
\selectlanguage{english}
\lettrine[lines=2]{I}{n} Thebais, the birthday of St. Paul, the first hermit who lived alone in 
 the desert from the sixteenth to the one hundred and thirteenth year of his 
 age. His soul was seen by St. Anthony carried by angels among the 
 choirs of apostles and prophets. His feast is kept on the 15th of this 
 month.
\switchcolumn*
\selectlanguage{latin}
In Cypro beáti Nicánoris, qui fuit unus de 
 septem primis Diáconis; atque, grátia fídei et virtúte admirándus, 
 gloriosíssime coronátus est.
\switchcolumn
\selectlanguage{english}
In Cyprus, blessed Nicanor, one of the first seven deacons, a man of 
 admirable faith and virtue, who received the crown of glory.
\switchcolumn*
\selectlanguage{latin}
Romæ sancti Agathónis Papæ, qui, sanctitáte et 
 doctrína conspícuus, quiévit in pace.
\switchcolumn
\selectlanguage{english}
At Rome, Pope St. Agatho, who, by a holy death, concluded a life remarkable 
 for sanctity and learning.
\switchcolumn*
\selectlanguage{latin}
Bitúricis, in Aquitánia, sancti Willhélmi, 
 Epíscopi et Confessóris, signis et virtútibus clari; quem Honórius Papa 
 Tértius in Sanctórum cánonem adscrípsit.
\switchcolumn
\selectlanguage{english}
At Bourges in Aquitaine, St. William, archbishop and confessor, renowned for 
 miracles and virtues. He was canonized by Pope Honorius III.
\switchcolumn*
\selectlanguage{latin}
Medioláni sancti Joánnis Boni, Epíscopi et Confessóris.
\switchcolumn
\selectlanguage{english}
At Milan, St. John the Good, bishop and confessor.
\switchcolumn*
\selectlanguage{latin}
Constantinópoli sancti Marciáni Presbyteri.
\switchcolumn
\selectlanguage{english}
At Constantinople, St. Marcian, priest.
\switchcolumn*
\selectlanguage{latin}
In monastério Cuxanénsi, in Gállia, natális sancti Petri Urséoli 
 Confessóris, qui, antea Venetiárum Dux et deínde Mónachus ex Ordine sancti 
 Benedícti, pietáte et virtútibus cláruit.
\switchcolumn
\selectlanguage{english}
In the monastery of Cusani in France, the birthday of St. Peter Orsini, 
 confessor, previously the Doge of Venice and afterwards monk of the Order of 
 St. Benedict, renowned for piety and miracles.
\switchcolumn*
\selectlanguage{latin}
Arétii, in Túscia, beáti Gregórii Décimi, civis Placentíni, qui, ex 
 Archidiácono Leodiénsi Summus Póntifex renuntiátus, Concílium Lugdunénse 
 secúndum celebrávit, Græcísque ad unitátem 
 fídei recéptis, compósitis Christianórum dissídiis, Terræ Sanctæ 
 recuperatióne institúta, de universáli Ecclésia, quam sanctíssime gubernávit, 
 óptime méritus est.
\switchcolumn
\selectlanguage{english}
At Arezzo in Tuscany, blessed Gregory X, a native of Piacenza, who was 
 elected Sovereign Pontiff while he was archdeacon of Liege. He held 
 the second Council of Lyons, received the Greeks into the unity of the 
 Church, appeased discords among the Christians, made generous efforts for 
 the recovery of the Holy Land, and governed the Church in a most holy 
 manner.
\switchcolumn*
\selectlanguage{latin}
\end{paracol}


% ---- martyrology/mart01/mart0111.htm
\needspace{10\baselineskip}
\begin{paracol}{2}
\selectlanguage{latin}
\begin{center}{\color{gregoriocolor} Tértio Idus Januárii. 
 Luna\dots\ }\end{center}
\switchcolumn
\selectlanguage{english}
\begin{center}{\color{gregoriocolor} The 11\textsuperscript{th} Day of 
 January. The\dots\ Day of the Moon.}\end{center}
\end{paracol}

\noindent\begin{tabularx}{\linewidth}{*{19}{>{\centering\arraybackslash}X}}
 \textcolor{gregoriocolor}{a} & \textcolor{gregoriocolor}{b} & \textcolor{gregoriocolor}{c} & \textcolor{gregoriocolor}{d} & \textcolor{gregoriocolor}{e} & \textcolor{gregoriocolor}{f} & \textcolor{gregoriocolor}{g} & \textcolor{gregoriocolor}{h} & \textcolor{gregoriocolor}{i} & \textcolor{gregoriocolor}{k} & \textcolor{gregoriocolor}{l} & \textcolor{gregoriocolor}{m} & \textcolor{gregoriocolor}{n} & \textcolor{gregoriocolor}{p} & \textcolor{gregoriocolor}{q} & \textcolor{gregoriocolor}{r} & \textcolor{gregoriocolor}{s} & \textcolor{gregoriocolor}{t} & \textcolor{gregoriocolor}{u} \\
 12 & 13 & 14 & 15 & 16 & 17 & 18 & 19 & 20 & 21 & 22 & 23 & 24 & 25 & 26 & 27 & 28 & 29 & 30 \\
\end{tabularx}
\vspace{0.5\baselineskip}
\noindent\begin{tabularx}{\linewidth}{*{12}{>{\centering\arraybackslash}X}}
 \textcolor{gregoriocolor}{A} & \textcolor{gregoriocolor}{B} & \textcolor{gregoriocolor}{C} & \textcolor{gregoriocolor}{D} & \textcolor{gregoriocolor}{E} & F & \textcolor{gregoriocolor}{F} & \textcolor{gregoriocolor}{G} & \textcolor{gregoriocolor}{H} & \textcolor{gregoriocolor}{M} & \textcolor{gregoriocolor}{N} & \textcolor{gregoriocolor}{P} \\
 1 & 2 & 3 & 4 & 5 & 6 & 6 & 7 & 8 & 9 & 10 & 11 \\
\end{tabularx}

\begin{paracol}{2}
\selectlanguage{latin}
\lettrine[lines=2]{R}{omæ} sancti Hygíni, Papæ et Mártyris; qui, in 
 persecutióne Antoníni, glorióse martyrium consummávit.
\switchcolumn
\selectlanguage{english}
\lettrine[lines=2]{A}{t} Rome, St. Hyginus, pope, who suffered a glorious martyrdom in the 
 persecution of Antoninus.
\switchcolumn*
\selectlanguage{latin}
Item Romæ natális sancti Melchíadis, Papæ et 
 Mártyris; qui multa, in persecutióne Maximiáni, passus est, atque, réddita 
 Ecclésiæ pace, quiévit in Dómino. Ipsíus autem festívitas quarto Idus 
 Decémbris celebrátur.
\switchcolumn
\selectlanguage{english}
Also at Rome, the birthday of St. Melchiades, who, having suffered much in 
 the persecution of Maximian, went to his rest in the Lord after peace 
 returned to the Church. His feast day is on the 10th of December.
\switchcolumn*
\selectlanguage{latin}
Firmi, in Picéno, sancti Alexándri, Epíscopi et Mártyris.
\switchcolumn
\selectlanguage{english}
At Fermo in Piceno, St. Alexander, bishop and martyr.
\switchcolumn*
\selectlanguage{latin}
Ambiáni, in Gállia, sancti Sálvii, Epíscopi et Mártyris.
\switchcolumn
\selectlanguage{english}
At Amiens in France, St. Salvius, bishop and martyr.
\switchcolumn*
\selectlanguage{latin}
In Africa beáti Sálvii Mártyris, in cujus natáli sanctus Augustínus sermónem 
 hábuit ad pópulum Carthaginénsem.
\switchcolumn
\selectlanguage{english}
In Africa, blessed Salvius, martyr, on whose birthday St. Augustine preached 
 to the people of Carthage.
\switchcolumn*
\selectlanguage{latin}
Alexandríæ sanctórum Mártyrum Petri, Sevéri et 
 Léucii.
\switchcolumn
\selectlanguage{english}
At Alexandria, the holy martyrs Peter, Severus and Leucius.
\switchcolumn*
\selectlanguage{latin}
Brundúsii sancti Léucii, Epíscopi et 
 Confessóris.
\switchcolumn
\selectlanguage{english}
At Brindisi, St. Leucius, bishop and confessor.
\switchcolumn*
\selectlanguage{latin}
In Judæa sancti Theodósii Cœnobiárchæ, in vico 
Cappadóciæ Magariásso nati; qui, multa passus pro fide cathólica, in pace 
tandem quiévit in eo monastério, quod ille super solitárium Hierosolymitánæ 
diœcésis montem exstrúxerat.
\switchcolumn
\selectlanguage{english}
In Judea, St. Theodosius, abbot, born in Cappadocia in the village of 
 Magarisso, who, after having endured great sufferings for the Catholic 
 faith, took his rest in peace at the monastery which he had erected on a 
 lonely hill in the diocese of Jerusalem.
\switchcolumn*
\selectlanguage{latin}
In Thebáide sancti Palǽmonis Abbátis, qui fuit 
 magíster sancti Pachómii.
\switchcolumn
\selectlanguage{english}
In Thebais, St. Palaemon, abbot, who was the teacher of St. Pachomius.
\switchcolumn*
\selectlanguage{latin}
Suppentóniæ, apud montem Soráctem, sancti 
 Anastásii Mónachi, et Sociórum; qui, divínitus vocáti, felíciter migravérunt 
 ad Dóminum.
\switchcolumn
\selectlanguage{english}
At Suppentonia, near Mount Soracte, St. Athanasius, monk, and his 
 companions, who were called by a voice from heaven to enter the kingdom of 
 God.
\switchcolumn*
\selectlanguage{latin}
Papíæ sanctæ Honorátæ Vírginis.
\switchcolumn
\selectlanguage{english}
At Pavia, St. Honorata, virgin.
\switchcolumn*
\selectlanguage{latin}
\end{paracol}


% ---- martyrology/mart01/mart0112.htm
\needspace{10\baselineskip}
\begin{paracol}{2}
\selectlanguage{latin}
\begin{center}{\color{gregoriocolor} Prídie Idus Januárii. 
 Luna\dots\ }\end{center}
\switchcolumn
\selectlanguage{english}
\begin{center}{\color{gregoriocolor} The 12\textsuperscript{th} Day of 
 January. The\dots\ Day of the Moon.}\end{center}
\end{paracol}

\noindent\begin{tabularx}{\linewidth}{*{19}{>{\centering\arraybackslash}X}}
 \textcolor{gregoriocolor}{a} & \textcolor{gregoriocolor}{b} & \textcolor{gregoriocolor}{c} & \textcolor{gregoriocolor}{d} & \textcolor{gregoriocolor}{e} & \textcolor{gregoriocolor}{f} & \textcolor{gregoriocolor}{g} & \textcolor{gregoriocolor}{h} & \textcolor{gregoriocolor}{i} & \textcolor{gregoriocolor}{k} & \textcolor{gregoriocolor}{l} & \textcolor{gregoriocolor}{m} & \textcolor{gregoriocolor}{n} & \textcolor{gregoriocolor}{p} & \textcolor{gregoriocolor}{q} & \textcolor{gregoriocolor}{r} & \textcolor{gregoriocolor}{s} & \textcolor{gregoriocolor}{t} & \textcolor{gregoriocolor}{u} \\
 13 & 14 & 15 & 16 & 17 & 18 & 19 & 20 & 21 & 22 & 23 & 24 & 25 & 26 & 27 & 28 & 29 & 30 & 1 \\
\end{tabularx}
\vspace{0.5\baselineskip}
\noindent\begin{tabularx}{\linewidth}{*{12}{>{\centering\arraybackslash}X}}
 \textcolor{gregoriocolor}{A} & \textcolor{gregoriocolor}{B} & \textcolor{gregoriocolor}{C} & \textcolor{gregoriocolor}{D} & \textcolor{gregoriocolor}{E} & F & \textcolor{gregoriocolor}{F} & \textcolor{gregoriocolor}{G} & \textcolor{gregoriocolor}{H} & \textcolor{gregoriocolor}{M} & \textcolor{gregoriocolor}{N} & \textcolor{gregoriocolor}{P} \\
 2 & 3 & 4 & 5 & 6 & 7 & 7 & 8 & 9 & 10 & 11 & 12 \\
\end{tabularx}

\begin{paracol}{2}
\selectlanguage{latin}
\lettrine[lines=2]{R}{omæ} sanctæ Tatiánæ Mártyris, quæ, sub 
 Alexándro Imperatóre, uncis atque pectínibus laniáta, béstiis expósita et in 
 ignem missa, sed nil læsa, demum, gládio percússa, migrávit in cælum.
\switchcolumn
\selectlanguage{english}
\lettrine[lines=2]{A}{t} Rome, in the time of Emperor Alexander, St. Tatiana, martyr, who had her 
 flesh torn with iron hooks and combs, was thrown to the beasts and cast into 
 the fire, but, having received no injury, was beheaded, and thus went to 
 heaven.
\switchcolumn*
\selectlanguage{latin}
Constantinópoli sanctórum Tígrii Presbyteri, et Eutrópii Lectóris; qui, 
 Arcádii Imperatóris témpore, cum de incéndio quo Ecclésia princeps et 
 Senátus cúria conflagráverant, tamquam 
 per eos ad exsílium sancti Joánnis Chrysóstomi ulciscéndum excitáto, per 
 calúmniam accusáti essent, sub Præfécto urbis 
 Optáto, inánium deórum superstitióne implícito et Christiánæ religiónis 
 osóre, passi sunt.
\switchcolumn
\selectlanguage{english}
At Constantinople, the Saints Tigrius, priest, and Eutropius, lector, who, 
 in the time of Emperor Arcadius, were falsely accused of the fire which 
 destroyed the principal church and the senate building in order to avenge 
 the exile of St. John Chrysostom. They suffered under Optatus, prefect 
 of the city, who was given to the worship of false gods and a hatred for the 
 Christian religion.
\switchcolumn*
\selectlanguage{latin}
In Achája sancti Sátyri Mártyris, qui, cum ante 
 quoddam idólum transíret, in illud exsufflávit, signans sibi frontem, atque 
 statim idólum córruit; ob quam causam decollátus est.
\switchcolumn
\selectlanguage{english}
In Achaia, St. Satyrus, martyr. As he passed before an idol and 
 breathed upon it, making the sign of the cross upon his forehead, the idol 
 immediately fell to the ground; for this reason he was beheaded.
\switchcolumn*
\selectlanguage{latin}
Eódem die sancti Arcádii Mártyris, génere et miráculis clari.
\switchcolumn
\selectlanguage{english}
On the same day, St. Arcadius, martyr, illustrious for his noble extraction 
 and miracles.
\switchcolumn*
\selectlanguage{latin}
In Africa sanctórum Mártyrum Zótici, Rogáti, 
 Modésti, Cástuli, et corónæ mílitum quadragínta.
\switchcolumn
\selectlanguage{english}
In Africa, the holy martyrs Zoticus, Rogatus, Modestus, Castulus, and forty 
 soldiers gloriously crowned.
\switchcolumn*
\selectlanguage{latin}
Tíbure sancti Zótici Mártyris.
\switchcolumn
\selectlanguage{english}
At Tivoli, St. Zoticus, martyr.
\switchcolumn*
\selectlanguage{latin}
Ephesi pássio sanctórum quadragínta duórum Monachórum, qui ob cultum 
 sanctárum Imáginum, sub Constantíno Coprónymo, 
 sævíssime cruciáti, martyrium complevérunt.
\switchcolumn
\selectlanguage{english}
At Ephesus, under Constantine Copronymus, the passion of forty-two holy 
 monks, who endured martyrdom after being most cruelly tortured for the 
 defence of sacred images.
\switchcolumn*
\selectlanguage{latin}
Ravénnæ sancti Joánnis, Epíscopi et Confessóris.
\switchcolumn
\selectlanguage{english}
At Ravenna, St. John, bishop and confessor.
\switchcolumn*
\selectlanguage{latin}
Verónæ sancti Probi Epíscopi.
\switchcolumn
\selectlanguage{english}
At Verona, St. Probus, bishop.
\switchcolumn*
\selectlanguage{latin}
In Anglia sancti Benedícti, Abbátis et Confessóris.
\switchcolumn
\selectlanguage{english}
In England, St. Benedict Biscop, abbot and confessor.
\switchcolumn*
\selectlanguage{latin}
\end{paracol}


% ---- martyrology/mart01/mart0113.htm
\needspace{10\baselineskip}
\begin{paracol}{2}
\selectlanguage{latin}
\begin{center}{\color{gregoriocolor} Idibus Januárii. 
 Luna\dots\ }\end{center}
\switchcolumn
\selectlanguage{english}
\begin{center}{\color{gregoriocolor} The 13\textsuperscript{th} Day of 
 January. The\dots\ Day of the Moon.}\end{center}
\end{paracol}

\noindent\begin{tabularx}{\linewidth}{*{19}{>{\centering\arraybackslash}X}}
 \textcolor{gregoriocolor}{a} & \textcolor{gregoriocolor}{b} & \textcolor{gregoriocolor}{c} & \textcolor{gregoriocolor}{d} & \textcolor{gregoriocolor}{e} & \textcolor{gregoriocolor}{f} & \textcolor{gregoriocolor}{g} & \textcolor{gregoriocolor}{h} & \textcolor{gregoriocolor}{i} & \textcolor{gregoriocolor}{k} & \textcolor{gregoriocolor}{l} & \textcolor{gregoriocolor}{m} & \textcolor{gregoriocolor}{n} & \textcolor{gregoriocolor}{p} & \textcolor{gregoriocolor}{q} & \textcolor{gregoriocolor}{r} & \textcolor{gregoriocolor}{s} & \textcolor{gregoriocolor}{t} & \textcolor{gregoriocolor}{u} \\
 14 & 15 & 16 & 17 & 18 & 19 & 20 & 21 & 22 & 23 & 24 & 25 & 26 & 27 & 28 & 29 & 30 & 1 & 2 \\
\end{tabularx}
\vspace{0.5\baselineskip}
\noindent\begin{tabularx}{\linewidth}{*{12}{>{\centering\arraybackslash}X}}
 \textcolor{gregoriocolor}{A} & \textcolor{gregoriocolor}{B} & \textcolor{gregoriocolor}{C} & \textcolor{gregoriocolor}{D} & \textcolor{gregoriocolor}{E} & F & \textcolor{gregoriocolor}{F} & \textcolor{gregoriocolor}{G} & \textcolor{gregoriocolor}{H} & \textcolor{gregoriocolor}{M} & \textcolor{gregoriocolor}{N} & \textcolor{gregoriocolor}{P} \\
 3 & 4 & 5 & 6 & 7 & 8 & 8 & 9 & 10 & 11 & 12 & 13 \\
\end{tabularx}

\begin{paracol}{2}
\selectlanguage{latin}
\lettrine[lines=1]{O}{ctáva} Epiphaníæ Domini.
\switchcolumn
\selectlanguage{english}
\lettrine[lines=1]{T}{he} Octave of the Epiphany of our Lord.
\switchcolumn*
\selectlanguage{latin}
Pictávis, in Gállia, natális sancti Hilárii, Epíscopi et Confessóris; qui, 
 ob cathólicam fidem, quam strénue propugnávit, quadriénnio apud Phrygiam 
 relegátus, ibi, inter ália mirácula, mórtuum suscitávit. Eum Pius 
 Nonus, Póntifex Máximus, universális Ecclésiæ 
 Doctórem declarávit et confirmávit. Ipsíus autem festum sequénti die 
 celebrátur.
\switchcolumn
\selectlanguage{english}
At Poitiers in France, the birthday of St. Hilary, bishop and confessor of 
 the Catholic faith which he courageously defended, and for which he was 
 banished for four years to Phrygia, where, among other miracles, he raised a 
 man from the dead. Pius IX declared him a doctor of the Church. 
 His festival is celebrated tomorrow.
\switchcolumn*
\selectlanguage{latin}
Rhemis, in Gállia, item natális sancti Remígii, Epíscopi et Confessóris. 
 Hic gentem Francórum convértit ad Christum, Clodovéo, 
 ipsórum Rege, sacris baptísmatis undis et fídei sacraméntis initiáto; et, 
 cum annos plúrimos in Episcopátu explésset, sanctitáte et miraculórum glória 
 conspícuus, decéssit e vita. Ejus vero festívitas Kaléndis Octóbris 
 potíssimum recólitur, quo die sacrum ipsíus corpus translátum fuit.
\switchcolumn
\selectlanguage{english}
At Rheims in France, St. Remigius, bishop and and confessor, who converted 
 the Franks to Christ, and brought Clovis, their king, to the holy font of 
 baptism and instructed him in the mysteries of faith. After he had 
 been bishop for many years, and had distinguished himself by his sanctity 
 and the power of working miracles, he departed this life. His feast is 
 kept on the 1st of October, on which day his holy body was transferred.
\switchcolumn*
\selectlanguage{latin}
Romæ, via Lavicána, corónæ sanctórum mílitum 
 quadragínta, quas ipsi, sub Galliéno Imperatóre, pro veræ fídei confessióne 
 percípere meruérunt.
\switchcolumn
\selectlanguage{english}
At Rome, on the Via Lavicana, the crowning of forty holy soldiers, a reward 
 they merited by confessing the true faith under Emperor Gallienus.
\switchcolumn*
\selectlanguage{latin}
Córdubæ, in Hispánia, sanctórum Mártyrum 
 Gumesíndi Presbyteri, et Servidéi Mónachi.
\switchcolumn
\selectlanguage{english}
At Cordova, the holy martyrs Gumesind, priest, and Servideus, monk.
\switchcolumn*
\selectlanguage{latin}
In Sardínia sancti Potíti Mártyris, qui, sub Antoníno Imperatóre et Gelásio 
 Præside, multa passus, demum gládio martyrium 
 consecútus est.
\switchcolumn
\selectlanguage{english}
In Sardinia, St. Potitus, martyr, who, having suffered much under Emperor 
 Antoninus and the governor Gelasius, was at last put to death by the sword.
\switchcolumn*
\selectlanguage{latin}
Singidóni, in Mysia superióre, sanctórum 
 Mártyrum Hérmyli et Stratoníci, qui, post sæva torménta, sub Licínio 
 Imperatóre, in Istrum flumen demérsi sunt.
\switchcolumn
\selectlanguage{english}
At Belgrade in Serbia, the holy martyrs Hermylus and Stratonicus, who were 
 severely tormented under Emperor Licinius, and then drowned in the river 
 Danube.
\switchcolumn*
\selectlanguage{latin}
Cæsaréæ, in Cappadócia, sancti Leóntii Epíscopi, 
 qui, sub Licínio, advérsus Gentíles, et, sub Constantíno, advérsus Ariános 
 plúrimum decertávit.
\switchcolumn
\selectlanguage{english}
At Caesarea in Cappadocia, St. Leontius, bishop, who fought strongly against 
 the heathens in the reign of Licinius, and against the Arians in the reign 
 of Constantine.
\switchcolumn*
\selectlanguage{latin}
Tréviris sancti Agrítii Epíscopi.
\switchcolumn
\selectlanguage{english}
At Treves, St. Agritius, bishop.
\switchcolumn*
\selectlanguage{latin}
In Versíaco monastério, in Gállia, sancti 
 Vivéntii Confessóris.
\switchcolumn
\selectlanguage{english}
In the monastery of Verzy in France, St. Viventius, confessor.
\switchcolumn*
\selectlanguage{latin}
Amaséæ, in Ponto, sanctæ Gláphyræ Vírginis.
\switchcolumn
\selectlanguage{english}
At Amasea in Pontus, St. Glaphyra, virgin.
\switchcolumn*
\selectlanguage{latin}
Medioláni, in cœnóbio sanctæ Marthæ, Beátæ 
 Verónicæ de Binásco Vírginis, ex Ordine sancti Augustíni.
\switchcolumn
\selectlanguage{english}
At Milan, in the monastery of St. Martha, blessed Veronica of Binasco, 
 virgin, of the Order of St. Augustine.
\switchcolumn*
\selectlanguage{latin}
\end{paracol}


% ---- martyrology/mart01/mart0114.htm
\needspace{10\baselineskip}
\begin{paracol}{2}
\selectlanguage{latin}
\begin{center}{\color{gregoriocolor} Décimo nono Kaléndas Februárii. 
 Luna\dots\ }\end{center}
\switchcolumn
\selectlanguage{english}
\begin{center}{\color{gregoriocolor} The 14\textsuperscript{th} Day of 
 January. The\dots\ Day of the Moon.}\end{center}
\end{paracol}

\noindent\begin{tabularx}{\linewidth}{*{19}{>{\centering\arraybackslash}X}}
 \textcolor{gregoriocolor}{a} & \textcolor{gregoriocolor}{b} & \textcolor{gregoriocolor}{c} & \textcolor{gregoriocolor}{d} & \textcolor{gregoriocolor}{e} & \textcolor{gregoriocolor}{f} & \textcolor{gregoriocolor}{g} & \textcolor{gregoriocolor}{h} & \textcolor{gregoriocolor}{i} & \textcolor{gregoriocolor}{k} & \textcolor{gregoriocolor}{l} & \textcolor{gregoriocolor}{m} & \textcolor{gregoriocolor}{n} & \textcolor{gregoriocolor}{p} & \textcolor{gregoriocolor}{q} & \textcolor{gregoriocolor}{r} & \textcolor{gregoriocolor}{s} & \textcolor{gregoriocolor}{t} & \textcolor{gregoriocolor}{u} \\
 15 & 16 & 17 & 18 & 19 & 20 & 21 & 22 & 23 & 24 & 25 & 26 & 27 & 28 & 29 & 30 & 1 & 2 & 3 \\
\end{tabularx}
\vspace{0.5\baselineskip}
\noindent\begin{tabularx}{\linewidth}{*{12}{>{\centering\arraybackslash}X}}
 \textcolor{gregoriocolor}{A} & \textcolor{gregoriocolor}{B} & \textcolor{gregoriocolor}{C} & \textcolor{gregoriocolor}{D} & \textcolor{gregoriocolor}{E} & F & \textcolor{gregoriocolor}{F} & \textcolor{gregoriocolor}{G} & \textcolor{gregoriocolor}{H} & \textcolor{gregoriocolor}{M} & \textcolor{gregoriocolor}{N} & \textcolor{gregoriocolor}{P} \\
 4 & 5 & 6 & 7 & 8 & 9 & 9 & 10 & 11 & 12 & 13 & 14 \\
\end{tabularx}

\begin{paracol}{2}
\selectlanguage{latin}
\lettrine[lines=2]{S}{ancti} Hilárii, Epíscopi Pictaviénsis, Confessóris et Ecclésiæ 
 Doctóris; qui prídie hujus diéi evolávit in cælum.
\switchcolumn
\selectlanguage{english}
\lettrine[lines=2]{S}{t.} Hilary, bishop of Poitiers, confessor and doctor of the Church, who 
 entered heaven on the thirteenth day of this month.
\switchcolumn*
\selectlanguage{latin}
Nolæ, in Campánia, natális sancti Felícis 
 Presbyteri, qui (ut sanctus Paulínus Epíscopus scribit), cum a 
 persecutóribus post torménta in cárcerem missus esset, et cóchleis ac 
 téstulis vinctus superpósitus jacéret, nocte ab Angelo solútus atque edúctus 
 fuit; póstmodum vero, cessánte persecutióne, ibídem, cum multos ad Christi 
 fidem exémplo vitæ ac doctrína convertísset, clarus miráculis quiévit in 
 pace.
\switchcolumn
\selectlanguage{english}
At Nola in Campania, the birthday of St. Felix, priest, who (as is related 
 by bishop St. Paulinus), after being subjected to torments by the 
 persecutors, was cast into prison, bound hand and foot, and extended on 
 shells and broken earthenware. In the night, however, his bonds were 
 loosened and he was delivered by an angel. The persecution over, he 
 brought many to the faith of Christ by his exemplary life and teaching, and, 
 renowned for miracles, rested in peace.
\switchcolumn*
\selectlanguage{latin}
In Judæa sancti Malachíæ Prophétæ.
\switchcolumn
\selectlanguage{english}
In Judea, St. Malachy, prophet.
\switchcolumn*
\selectlanguage{latin}
In monte Sina sanctórum trigínta octo Monachórum, a Saracénis ob Christi 
 fidem interfectórum.
\switchcolumn
\selectlanguage{english}
On Mount Sinai, thirty-eight holy monks killed by the Saracens for the faith 
 of Christ.
\switchcolumn*
\selectlanguage{latin}
In Rhaíthi regióne, in Ægypto, sanctórum 
 quadragínta trium Monachórum, qui, pro Christiána religióne, a Blémmiis 
 occísi sunt.
\switchcolumn
\selectlanguage{english}
In Egypt, in the district of Raithy, forty-three holy monks, who were put to 
 death by the Blemmians for the Christian religion.
\switchcolumn*
\selectlanguage{latin}
Medioláni sancti Dátii, Epíscopi et Confessóris; cujus méminit beátus 
 Gregórius Papa.
\switchcolumn
\selectlanguage{english}
At Milan, St. Datius, bishop and confessor, mentioned by pope St. Gregory.
\switchcolumn*
\selectlanguage{latin}
In Africa sancti Euphrásii Epíscopi.
\switchcolumn
\selectlanguage{english}
In Africa, St. Euphrasius, bishop.
\switchcolumn*
\selectlanguage{latin}
Neocæsaréæ, in Ponto, sanctæ Macrínæ, discípulæ 
 beáti Gregórii Thaumatúrgi, et áviæ sancti Basilíi, quæ eúndem Basilíum 
 educávit in fide.
\switchcolumn
\selectlanguage{english}
At Neocaesarea in Pontus, St. Macrina, disciple of St. Gregory the 
 Wonder-Worker, and grandmother of St. Basil, whom she educated in the 
 Christian faith.
\switchcolumn*
\selectlanguage{latin}
\end{paracol}


% ---- martyrology/mart01/mart0115.htm
\needspace{10\baselineskip}
\begin{paracol}{2}
\selectlanguage{latin}
\begin{center}{\color{gregoriocolor} Décimo octávo Kaléndas Februárii. 
 Luna\dots\ }\end{center}
\switchcolumn
\selectlanguage{english}
\begin{center}{\color{gregoriocolor} The 15\textsuperscript{th} Day of 
 January. The\dots\ Day of the Moon.}\end{center}
\end{paracol}

\noindent\begin{tabularx}{\linewidth}{*{19}{>{\centering\arraybackslash}X}}
 \textcolor{gregoriocolor}{a} & \textcolor{gregoriocolor}{b} & \textcolor{gregoriocolor}{c} & \textcolor{gregoriocolor}{d} & \textcolor{gregoriocolor}{e} & \textcolor{gregoriocolor}{f} & \textcolor{gregoriocolor}{g} & \textcolor{gregoriocolor}{h} & \textcolor{gregoriocolor}{i} & \textcolor{gregoriocolor}{k} & \textcolor{gregoriocolor}{l} & \textcolor{gregoriocolor}{m} & \textcolor{gregoriocolor}{n} & \textcolor{gregoriocolor}{p} & \textcolor{gregoriocolor}{q} & \textcolor{gregoriocolor}{r} & \textcolor{gregoriocolor}{s} & \textcolor{gregoriocolor}{t} & \textcolor{gregoriocolor}{u} \\
 16 & 17 & 18 & 19 & 20 & 21 & 22 & 23 & 24 & 25 & 26 & 27 & 28 & 29 & 30 & 1 & 2 & 3 & 4 \\
\end{tabularx}
\vspace{0.5\baselineskip}
\noindent\begin{tabularx}{\linewidth}{*{12}{>{\centering\arraybackslash}X}}
 \textcolor{gregoriocolor}{A} & \textcolor{gregoriocolor}{B} & \textcolor{gregoriocolor}{C} & \textcolor{gregoriocolor}{D} & \textcolor{gregoriocolor}{E} & F & \textcolor{gregoriocolor}{F} & \textcolor{gregoriocolor}{G} & \textcolor{gregoriocolor}{H} & \textcolor{gregoriocolor}{M} & \textcolor{gregoriocolor}{N} & \textcolor{gregoriocolor}{P} \\
 5 & 6 & 7 & 8 & 9 & 10 & 10 & 11 & 12 & 13 & 14 & 15 \\
\end{tabularx}

\begin{paracol}{2}
\selectlanguage{latin}
\lettrine[lines=2]{S}{ancti} Pauli, primi Eremítæ, Confessóris; qui 
 quarto Idus Januárii inter beatórum ágmina translátus fuit.
\switchcolumn
\selectlanguage{english}
\lettrine[lines=2]{S}{t.} Paul, the first hermit, who was carried to the home of the blessed on 
 the tenth of this month.
\switchcolumn*
\selectlanguage{latin}
In território Andegavénsi beáti Mauri Abbátis, qui fuit discípulus sancti 
 Benedícti; et, hujus disciplínis usque ab infántia erudítus, quantum in eis 
 profécerit, inter ália quæ apud eum pósitus 
 gessit (res nova et post Petrum fere inusitáti), pédibus super aquas 
 incédens patefécit. In Gállias inde ab ipso Benedícto diréctus, ibi, 
 constrúcto célebri monastério, cui quadragínta annis præfuit, miraculórum 
 glória clarus, in pace quiévit.
\switchcolumn
\selectlanguage{english}
In the diocese of Angers, blessed Maurus, abbot and disciple of St. Benedict. 
 Beginning his discipline in infancy, he made great progress with so able a 
 master, for while he was still under the saint's instruction he miraculously 
 walked upon the water, a prodigy unheard of since the days of St. Peter. 
 Sent later to France by St. Benedict, he built a famous monastery, which he 
 governed for forty years, and after performing striking miracles, he rested 
 in peace.
\switchcolumn*
\selectlanguage{latin}
In Judæa sanctórum Hábacuc et Michǽæ 
 Prophetárum, quorum corpora, sub Theodósio senióre, divína revelatióne sunt 
 repérta.
\switchcolumn
\selectlanguage{english}
In Judea, the holy prophets Habakkuk and Micah, whose bodies were found by 
 divine revelation in the days of Theodosius the Elder.
\switchcolumn*
\selectlanguage{latin}
Cárali, in Sardínia, sancti Ephísii Mártyris, 
 qui, in persecutióne Diocletiáni, sub Flaviáno Júdice, plúrimis torméntis 
 divína virtúte superátis, demum, abscíssis cervícibus, victor migrávit in 
 cælum.
\switchcolumn
\selectlanguage{english}
At Cagliari in Sardinia, St. Ephisius, martyr, who, in the persecution of 
 Diocletian and under the judge Flavian, having, by the assistance of God, 
 overcome many torments, was beheaded and ascended to heaven.
\switchcolumn*
\selectlanguage{latin}
Anágniæ 
 sanctæ Secundínæ, Vírginis et Mártyris; quæ sub Décio Imperatóre passa est.
\switchcolumn
\selectlanguage{english}
At Anagni, St. Secundina, virgin and martyr, who suffered under Emperor 
 Decius.
\switchcolumn*
\selectlanguage{latin}
Nolæ, in Campánia, sancti Máximi Epíscopi.
\switchcolumn
\selectlanguage{english}
At Nola in Campania, St. Maximus, bishop.
\switchcolumn*
\selectlanguage{latin}
Arvérnis, in Gállia, sancti Boníti, Epíscopi et Confessóris.
\switchcolumn
\selectlanguage{english}
In Auvergne in France, St. Bonitus, bishop and confessor.
\switchcolumn*
\selectlanguage{latin}
In Ægypto sancti Macárii Abbátis, qui fuit 
 discípulus beáti Antónii, ac vita et miráculis celebérrimus éxstitit.
\switchcolumn
\selectlanguage{english}
In Egypt, St. Macarius, abbot, disciple of St. Anthony, very celebrated for 
 his life and miracles.
\switchcolumn*
\selectlanguage{latin}
Alexandríæ beáti Isidóri, sanctitáte vitæ, fide 
 et miráculis clari.
\switchcolumn
\selectlanguage{english}
At Alexandria, blessed Isidore, renowned for holiness of life, faith, and 
 miracles.
\switchcolumn*
\selectlanguage{latin}
Constantinópoli sancti Joánnis Calybítæ, qui 
 aliquándiu in ángulo domus patérnæ, deínde in tugúrio, ignótus paréntibus, 
 habitávit; a quibus in morte ágnitus, miráculis cláruit. Ipsíus corpus 
 póstea Romam translátum, et in Insulæ Tiberínæ Ecclésia, in ejus honórem 
 erécta, collocátum est.
\switchcolumn
\selectlanguage{english}
At Constantinople, St. John Calybita. For some time living unknown to 
 his parents in a corner of their house, and later in a hut on an island in 
 the Tiber, he was recognized by them only at his death. Being renowned 
 for miracles, his body was afterwards taken to Rome and buried on the Island 
 in the Tiber, where a church was subsequently erected in his honour.
\switchcolumn*
\selectlanguage{latin}
\end{paracol}


% ---- martyrology/mart01/mart0116.htm
\needspace{10\baselineskip}
\begin{paracol}{2}
\selectlanguage{latin}
\begin{center}{\color{gregoriocolor} Décimo séptimo Kaléndas Februárii. 
 Luna\dots\ }\end{center}
\switchcolumn
\selectlanguage{english}
\begin{center}{\color{gregoriocolor} The 16\textsuperscript{th} Day of 
 January. The\dots\ Day of the Moon.}\end{center}
\end{paracol}

\noindent\begin{tabularx}{\linewidth}{*{19}{>{\centering\arraybackslash}X}}
 \textcolor{gregoriocolor}{a} & \textcolor{gregoriocolor}{b} & \textcolor{gregoriocolor}{c} & \textcolor{gregoriocolor}{d} & \textcolor{gregoriocolor}{e} & \textcolor{gregoriocolor}{f} & \textcolor{gregoriocolor}{g} & \textcolor{gregoriocolor}{h} & \textcolor{gregoriocolor}{i} & \textcolor{gregoriocolor}{k} & \textcolor{gregoriocolor}{l} & \textcolor{gregoriocolor}{m} & \textcolor{gregoriocolor}{n} & \textcolor{gregoriocolor}{p} & \textcolor{gregoriocolor}{q} & \textcolor{gregoriocolor}{r} & \textcolor{gregoriocolor}{s} & \textcolor{gregoriocolor}{t} & \textcolor{gregoriocolor}{u} \\
 17 & 18 & 19 & 20 & 21 & 22 & 23 & 24 & 25 & 26 & 27 & 28 & 29 & 30 & 1 & 2 & 3 & 4 & 5 \\
\end{tabularx}
\vspace{0.5\baselineskip}
\noindent\begin{tabularx}{\linewidth}{*{12}{>{\centering\arraybackslash}X}}
 \textcolor{gregoriocolor}{A} & \textcolor{gregoriocolor}{B} & \textcolor{gregoriocolor}{C} & \textcolor{gregoriocolor}{D} & \textcolor{gregoriocolor}{E} & F & \textcolor{gregoriocolor}{F} & \textcolor{gregoriocolor}{G} & \textcolor{gregoriocolor}{H} & \textcolor{gregoriocolor}{M} & \textcolor{gregoriocolor}{N} & \textcolor{gregoriocolor}{P} \\
 6 & 7 & 8 & 9 & 10 & 11 & 11 & 12 & 13 & 14 & 15 & 16 \\
\end{tabularx}

\begin{paracol}{2}
\selectlanguage{latin}
\lettrine[lines=2]{R}{omæ,} via Salária, natális sancti Marcélli 
 Primi, Papæ et Mártyris; qui, ob cathólicæ fídei confessiónem, jubénte 
 Maxéntio tyránno, primo cæsus est fústibus, deínde ad servítium animálium 
 cum custódia pública deputátus, et ibídem, serviéndo indútus amíctu cilícino, 
 defúnctus est.
\switchcolumn
\selectlanguage{english}
\lettrine[lines=2]{A}{t} Rome, on the Salarian Way, the birthday of Pope St. Marcellus I, a martyr 
 for the confession of the Catholic faith. By command of the tyrant 
 Maxentius he was beaten with clubs, then sent to take care of animals, with 
 a guard to watch him. In this servile office, dressed in haircloth, he 
 departed this life.
\switchcolumn*
\selectlanguage{latin}
Marróchii, in Africa, pássio sanctórum quinque Protomártyrum Ordinis Minórum, 
 scílicet Berárdi, Petri atque Othónis Sacerdótum, Accúrsii et Adjúti 
 Laicórum; qui, ob Christiánæ fídei 
 prædicatiónem ac Mahuméticæ reprobatiónem legis, post vária torménta et 
 ludíbria, a Saracenórum Rege, scissis gládio capítibus, enecáti sunt.
\switchcolumn
\selectlanguage{english}
In Morocco in Africa, the martyrdom of the five Protomartyrs of the Order of 
 Friars Minor, Berard, Peter, and Otto who were priests, and Accursius and 
 Adjutus who were lay brothers. For preaching the Catholic faith, and 
 because of their hatred of the Mohammedan Law, after various torments and 
 mockeries by the Saracen king, they were beheaded.
\switchcolumn*
\selectlanguage{latin}
Rhinocolúræ, in Ægypto, sancti Melæ Epíscopi, 
 qui, sub Valénte exsílium et ália grávia pro fide cathólica passus, in pace 
 quiévit.
\switchcolumn
\selectlanguage{english}
At Rhinocolura in Egypt, the holy bishop St. Melas, who rested in peace 
 after suffering exile and other painful trials for the Catholic faith during 
 the reign of Emperor Valens.
\switchcolumn*
\selectlanguage{latin}
Areláte, in Gállia, sancti Honoráti, Epíscopi et Confessóris; cujus vita tam 
 doctrína quam miráculis fuit illústris.
\switchcolumn
\selectlanguage{english}
At Arles in France, St. Honoratus, bishop and confessor, whose life was 
 renowned for learning and for miracles.
\switchcolumn*
\selectlanguage{latin}
Opitérgii, in Venetórum fínibus, sancti Titiáni, Epíscopi et Confessóris.
\switchcolumn
\selectlanguage{english}
At Oderzo near Venice, St. Titian, bishop and confessor.
\switchcolumn*
\selectlanguage{latin}
Fundis, in Látio, sancti Honoráti Abbátis, cujus méminit beátus Gregórius 
 Papa.
\switchcolumn
\selectlanguage{english}
At Fondi in Lazio, St. Honoratus, abbot, mentioned by Pope St. Gregory.
\switchcolumn*
\selectlanguage{latin}
In castro cui nomen Macériæ, ad Altéjam flúvium, 
 in Gállia, sancti Furséi Confessóris; cujus corpus ad monastérium Perónæ 
 póstmodum translátum est.
\switchcolumn
\selectlanguage{english}
At Froheins, in the diocese of Amiens in France, St. Fursey, confessor, 
 whose body was afterwards transferred to the monastery of Peronne.
\switchcolumn*
\selectlanguage{latin}
Romæ sanctæ Priscíllæ, quæ se súaque pio 
 Mártyrum obséquio mancipávit.
\switchcolumn
\selectlanguage{english}
At Rome, St. Priscilla, who devoted herself and her goods to the service of 
 the martyrs.
\switchcolumn*
\selectlanguage{latin}
\end{paracol}


% ---- martyrology/mart01/mart0117.htm
\needspace{10\baselineskip}
\begin{paracol}{2}
\selectlanguage{latin}
\begin{center}{\color{gregoriocolor} Décimo sexto Kaléndas Februárii. 
 Luna\dots\ }\end{center}
\switchcolumn
\selectlanguage{english}
\begin{center}{\color{gregoriocolor} The 17\textsuperscript{th} Day of 
 January. The\dots\ Day of the Moon.}\end{center}
\end{paracol}

\noindent\begin{tabularx}{\linewidth}{*{19}{>{\centering\arraybackslash}X}}
 \textcolor{gregoriocolor}{a} & \textcolor{gregoriocolor}{b} & \textcolor{gregoriocolor}{c} & \textcolor{gregoriocolor}{d} & \textcolor{gregoriocolor}{e} & \textcolor{gregoriocolor}{f} & \textcolor{gregoriocolor}{g} & \textcolor{gregoriocolor}{h} & \textcolor{gregoriocolor}{i} & \textcolor{gregoriocolor}{k} & \textcolor{gregoriocolor}{l} & \textcolor{gregoriocolor}{m} & \textcolor{gregoriocolor}{n} & \textcolor{gregoriocolor}{p} & \textcolor{gregoriocolor}{q} & \textcolor{gregoriocolor}{r} & \textcolor{gregoriocolor}{s} & \textcolor{gregoriocolor}{t} & \textcolor{gregoriocolor}{u} \\
 18 & 19 & 20 & 21 & 22 & 23 & 24 & 25 & 26 & 27 & 28 & 29 & 30 & 1 & 2 & 3 & 4 & 5 & 6 \\
\end{tabularx}
\vspace{0.5\baselineskip}
\noindent\begin{tabularx}{\linewidth}{*{12}{>{\centering\arraybackslash}X}}
 \textcolor{gregoriocolor}{A} & \textcolor{gregoriocolor}{B} & \textcolor{gregoriocolor}{C} & \textcolor{gregoriocolor}{D} & \textcolor{gregoriocolor}{E} & F & \textcolor{gregoriocolor}{F} & \textcolor{gregoriocolor}{G} & \textcolor{gregoriocolor}{H} & \textcolor{gregoriocolor}{M} & \textcolor{gregoriocolor}{N} & \textcolor{gregoriocolor}{P} \\
 7 & 8 & 9 & 10 & 11 & 12 & 12 & 13 & 14 & 15 & 16 & 17 \\
\end{tabularx}

\begin{paracol}{2}
\selectlanguage{latin}
\lettrine[lines=2]{I}{n} Thebáide sancti Antónii Abbátis, qui, multórum Monachórum Pater, vita et 
 miráculis præclaríssimus vixit; cujus gesta 
 sanctus Athanásius insígni volúmine prosecútus est. Ejus autem sacrum 
 corpus, sub Justiniáno Imperatóre, divína revelatióne repértum et 
 Alexandríam delátum, in Ecclésia sancti Joánnis Baptístæ humátum fuit.
\switchcolumn
\selectlanguage{english}
\lettrine[lines=2]{I}{n} Thebais, St. Anthony, abbot and spiritual guide of many monks, who was 
 most celebrated for his life and miracles of which St. Athanasius has 
 written a detailed account. His holy body was found by a divine 
 revelation during the reign of Emperor Justinian and brought to Alexandria, 
 where it was buried in the church of St. John Baptist.
\switchcolumn*
\selectlanguage{latin}
Apud Língonas, in Gállia, sanctórum 
 tergeminórum Speusíppi, Eleusíppi et Meleusíppi; qui, cum ávia sua Leonílla, 
 martyrio coronáti sunt, témpore Marci Aurélii Imperatóris.
\switchcolumn
\selectlanguage{english}
At Langres in France, in the time of Marcus Aurelius, the Saints Speusippus, 
 Eleusippus, and Meleusippus, born at one birth, were crowned with martyrdom 
 together with their grandmother Leonilla.
\switchcolumn*
\selectlanguage{latin}
Apud Bitúricas, in Aquitánia, deposítio sancti 
 Sulpícii Epíscopi, cognoménto Pii, cujus vita et mors pretiósa gloriósis 
 miráculis commendátur.
\switchcolumn
\selectlanguage{english}
At Bourges in Aquitaine, the death of the holy Bishop Sulpice, surnamed Pius, whose life 
 and precious death were approved by glorious miracles.
\switchcolumn*
\selectlanguage{latin}
Romæ, in monastério sancti Andréæ, beatórum 
 Monachórum Antónii, Méruli et Joánnis, de quibus scribit sanctus Gregórius 
 Papa.
\switchcolumn
\selectlanguage{english}
At Rome, in the monastery of St. Andrew, the blessed monks Anthony, Merulus, 
 and John, of whom Pope St. Gregory speaks in his writings.
\switchcolumn*
\selectlanguage{latin}
In fínibus Edessénæ regiónis, in Mesopotámia, sancti Juliáni Eremítæ, cognoménto Sabæ, qui, Valéntis Imperatóris témpore, 
 fidem cathólicam, Antiochíæ ferme collápsam, virtúte miraculórum eréxit.
\switchcolumn
\selectlanguage{english}
At Edessa in Mesopotamia, in the time of Emperor Valens, St. Julian Sabas 
 the Elder, who miraculously restored the Catholic faith at Antioch, although 
 it was almost destroyed in that city.
\switchcolumn*
\selectlanguage{latin}
Romæ Invéntio sanctórum Mártyrum Diodóri 
 Presbyteri, Mariáni Diáconi, et Sociórum; qui, sancto Stéphano Papa 
 Ecclésiam Dei regénte, martyrium Kaléndis Decémbris sunt assecúti.
\switchcolumn
\selectlanguage{english}
At Rome, the finding of the holy martyrs Diodorus, priest, and Marian, 
 deacon, and their companions. They suffered martyrdom on the 1st of 
 December during the pontificate of Pope St. Stephen.
\switchcolumn*
\selectlanguage{latin}
\end{paracol}


% ---- martyrology/mart01/mart0118.htm
\needspace{10\baselineskip}
\begin{paracol}{2}
\selectlanguage{latin}
\begin{center}{\color{gregoriocolor} Quintodécimo Kaléndas Februárii. 
 Luna\dots\ }\end{center}
\switchcolumn
\selectlanguage{english}
\begin{center}{\color{gregoriocolor} The 18\textsuperscript{th} Day of 
 January. The\dots\ Day of the Moon.}\end{center}
\end{paracol}

\noindent\begin{tabularx}{\linewidth}{*{19}{>{\centering\arraybackslash}X}}
 \textcolor{gregoriocolor}{a} & \textcolor{gregoriocolor}{b} & \textcolor{gregoriocolor}{c} & \textcolor{gregoriocolor}{d} & \textcolor{gregoriocolor}{e} & \textcolor{gregoriocolor}{f} & \textcolor{gregoriocolor}{g} & \textcolor{gregoriocolor}{h} & \textcolor{gregoriocolor}{i} & \textcolor{gregoriocolor}{k} & \textcolor{gregoriocolor}{l} & \textcolor{gregoriocolor}{m} & \textcolor{gregoriocolor}{n} & \textcolor{gregoriocolor}{p} & \textcolor{gregoriocolor}{q} & \textcolor{gregoriocolor}{r} & \textcolor{gregoriocolor}{s} & \textcolor{gregoriocolor}{t} & \textcolor{gregoriocolor}{u} \\
 19 & 20 & 21 & 22 & 23 & 24 & 25 & 26 & 27 & 28 & 29 & 30 & 1 & 2 & 3 & 4 & 5 & 6 & 7 \\
\end{tabularx}
\vspace{0.5\baselineskip}
\noindent\begin{tabularx}{\linewidth}{*{12}{>{\centering\arraybackslash}X}}
 \textcolor{gregoriocolor}{A} & \textcolor{gregoriocolor}{B} & \textcolor{gregoriocolor}{C} & \textcolor{gregoriocolor}{D} & \textcolor{gregoriocolor}{E} & F & \textcolor{gregoriocolor}{F} & \textcolor{gregoriocolor}{G} & \textcolor{gregoriocolor}{H} & \textcolor{gregoriocolor}{M} & \textcolor{gregoriocolor}{N} & \textcolor{gregoriocolor}{P} \\
 8 & 9 & 10 & 11 & 12 & 13 & 13 & 14 & 15 & 16 & 17 & 18 \\
\end{tabularx}

\begin{paracol}{2}
\selectlanguage{latin}
\lettrine[lines=2]{C}{áthedra} sancti Petri Apóstoli, qua primum Romæ 
 sedit.
\switchcolumn
\selectlanguage{english}
\lettrine[lines=2]{T}{he} Chair of St. Peter the Apostle, who established the Holy See at Rome.
\switchcolumn*
\selectlanguage{latin}
Ibídem pássio sanctæ Priscæ, Vírginis et 
 Mártyris; quæ sub Cláudio Imperatóre, post multa torménta, martyrio coronáta 
 est.
\switchcolumn
\selectlanguage{english}
In the same place, under Emperor Claudius, the passion of St. Prisca, virgin 
 and martyr, who, after undergoing many torments, was crowned with martyrdom.
\switchcolumn*
\selectlanguage{latin}
In Ponto natális sanctórum Mártyrum Moséi et 
 Ammónii, qui, cum essent mílites, primo ad metálla damnáti sunt, ac 
 novíssime igni tráditi.
\switchcolumn
\selectlanguage{english}
In Pontus, the birthday of the holy martyrs Mosseus and Ammonius, soldiers, 
 who were first condemned to work in the metal mines, then cast into the 
 fire.
\switchcolumn*
\selectlanguage{latin}
Ibídem sancti Athenógenis, antíqui Theólogi, 
 qui, per ignem consummatúrus martyrium, hymnum lætus cécinit, quem et 
 discípulis scriptum relíquit.
\switchcolumn
\selectlanguage{english}
In the same country, St. Athenogenes, an aged divine, who, on the point of 
 being martyred by fire, joyfully sang a hymn, which he left in writing to 
 his disciples.
\switchcolumn*
\selectlanguage{latin}
Turónis, in Gállia, sancti Volusiáni Epíscopi, qui, a Gothis captus, in 
 exsílio spíritum Deo réddidit.
\switchcolumn
\selectlanguage{english}
At Tours in France, St. Volusian, bishop, who was made captive by the Goths, 
 and in exile gave up his soul unto God.
\switchcolumn*
\selectlanguage{latin}
In monastério Lutrénsi, in Burgúndia, sancti Deícolæ 
 Abbátis, qui, natióne Hibérnus, discípulus fuit beáti Columbáni.
\switchcolumn
\selectlanguage{english}
In the monastery of Lure in Burgundy, St. Deicola, abbot, a native of 
 Ireland and a disciple of St. Columban.
\switchcolumn*
\selectlanguage{latin}
Turónis, in Gállia, sancti Leobárdi reclúsi, qui mira abstinéntia et 
 humilitáte refúlsit.
\switchcolumn
\selectlanguage{english}
At Tours in France, St. Leobard, anchoret, a man of wonderful abstinence and 
 humility.
\switchcolumn*
\selectlanguage{latin}
Novocómi sanctæ Liberátæ Vírginis.
\switchcolumn
\selectlanguage{english}
At Como, St. Liberata, virgin.
\switchcolumn*
\selectlanguage{latin}
Budæ, in Hungária, sanctæ Margarítæ Vírginis, e 
 régia Arpadénsium família, Ordinis sancti Domínici Moniális, virtúte 
 castitátis et arctíssima pæniténtia insígnis, quam Pius Duodécimus, Póntifex 
 Máximus, sanctárum Vírginum catálogo adscrípsit.
\switchcolumn
\selectlanguage{english}
At Buda in Hungary, St. Margaret, virgin, from the royal family of Arpad, 
 and a nun of the Order of St. Dominic, endued with the virtues of chastity 
 and a burning penitence. The Supreme Pontiff, Pius XII, added her to 
 the list of holy virgins.
\switchcolumn*
\selectlanguage{latin}
\end{paracol}


% ---- martyrology/mart01/mart0119.htm
\needspace{10\baselineskip}
\begin{paracol}{2}
\selectlanguage{latin}
\begin{center}{\color{gregoriocolor} Quartodécimo Kaléndas Februárii. 
 Luna\dots\ }\end{center}
\switchcolumn
\selectlanguage{english}
\begin{center}{\color{gregoriocolor} The 19\textsuperscript{th} Day of 
 January. The\dots\ Day of the Moon.}\end{center}
\end{paracol}

\noindent\begin{tabularx}{\linewidth}{*{19}{>{\centering\arraybackslash}X}}
 \textcolor{gregoriocolor}{a} & \textcolor{gregoriocolor}{b} & \textcolor{gregoriocolor}{c} & \textcolor{gregoriocolor}{d} & \textcolor{gregoriocolor}{e} & \textcolor{gregoriocolor}{f} & \textcolor{gregoriocolor}{g} & \textcolor{gregoriocolor}{h} & \textcolor{gregoriocolor}{i} & \textcolor{gregoriocolor}{k} & \textcolor{gregoriocolor}{l} & \textcolor{gregoriocolor}{m} & \textcolor{gregoriocolor}{n} & \textcolor{gregoriocolor}{p} & \textcolor{gregoriocolor}{q} & \textcolor{gregoriocolor}{r} & \textcolor{gregoriocolor}{s} & \textcolor{gregoriocolor}{t} & \textcolor{gregoriocolor}{u} \\
 20 & 21 & 22 & 23 & 24 & 25 & 26 & 27 & 28 & 29 & 30 & 1 & 2 & 3 & 4 & 5 & 6 & 7 & 8 \\
\end{tabularx}
\vspace{0.5\baselineskip}
\noindent\begin{tabularx}{\linewidth}{*{12}{>{\centering\arraybackslash}X}}
 \textcolor{gregoriocolor}{A} & \textcolor{gregoriocolor}{B} & \textcolor{gregoriocolor}{C} & \textcolor{gregoriocolor}{D} & \textcolor{gregoriocolor}{E} & F & \textcolor{gregoriocolor}{F} & \textcolor{gregoriocolor}{G} & \textcolor{gregoriocolor}{H} & \textcolor{gregoriocolor}{M} & \textcolor{gregoriocolor}{N} & \textcolor{gregoriocolor}{P} \\
 9 & 10 & 11 & 12 & 13 & 14 & 14 & 15 & 16 & 17 & 18 & 19 \\
\end{tabularx}

\begin{paracol}{2}
\selectlanguage{latin}
\lettrine[lines=2]{R}{omæ,} via Cornélia, sanctórum Mártyrum Márii et 
 Marthæ cónjugum, et filiórum Audífacis et Abachum, nobílium Persárum; qui 
 Romam, tempóribus Cláudii Príncipis, ad oratiónem vénerant. Ex eis 
 vero, post tolerátos fustes, equúleum, ignes, ungues férreos manuúmque 
 præcisiónem, Martha in Nympha necáta est; céteri sunt decolláti, et córpora 
 eórum incénsa.
\switchcolumn
\selectlanguage{english}
\lettrine[lines=2]{A}{t} Rome, on the Cornelian Road, the holy martyrs Marius and his wife Martha, 
 with their sons Audifax and Habbakuk, noble Persians, who came to Rome 
 through devotion in the time of Emperor Claudius. After they had been 
 beaten with rods, tormented on the rack and with fire, lacerated with iron 
 hooks, and had endured the cutting off of their hands, Martha was put to 
 death in the place called Nympha; the others were beheaded and cast into the 
 fire.
\switchcolumn*
\selectlanguage{latin}
Item sancti Canúti, Regis et Mártyris.
\switchcolumn
\selectlanguage{english}
Also St. Canute, king and martyr.
\switchcolumn*
\selectlanguage{latin}
Smyrnæ natális beáti Germánici Mártyris, qui, 
 sub Marco Antoníno et Lúcio Aurélio, cum primævæ ætátis venustáte floréret, 
 damnátus a Júdice, et, per grátiam virtútis Dei, metum corpóreæ fragilitátis 
 exclúdens, præparátum sibi béstiam sponte provocávit; cujus déntibus 
 comminútus, vero pani Dómino Jesu Christo, pro ipso móriens, méruit incorporári.
\switchcolumn
\selectlanguage{english}
At Smyrna, under Marcus Antoninus and Lucius Aurelius, the birthday of 
 blessed Germanicus, martyr, who, in the bloom of youth, being strengthened 
 by the grace of God, and banishing all fear, provoked the beast which, by 
 order of the judge, was to devour him. Being ground by its teeth, he 
 deserved to be incorporated into the true Bread of Life, Christ Jesus, for 
 whom he died.
\switchcolumn*
\selectlanguage{latin}
In Africa sanctórum Mártyrum Pauli, Geróntii, Januárii, Saturníni, Succéssi, 
 Júlii, Cati, Piæ et Germánæ.
\switchcolumn
\selectlanguage{english}
In Africa., the holy martyrs Paul, Gerontius, Januarius, Saturninus, 
 Successus, Julius, Catus, Pia, and Germana.
\switchcolumn*
\selectlanguage{latin}
Apud Spolétum pássio sancti Pontiáni Mártyris, qui, témpore Antoníni 
 Imperatóris, a Fabiáno Júdice, pro Christo vehementíssime virgis cæsus, 
 jussus est super carbónes nudis pédibus ambuláre, sed a carbónibus nil læsus, 
 equúleo et uncínis férreis jussus est suspéndi, et sic in cárcerem trudi, 
 ubi Angélica visitatióne méruit confortári; postque leónibus expósitus et 
 plumbo fervénti perfúsus, tandem gládio percússus est.
\switchcolumn
\selectlanguage{english}
At Spoleto, in the days of Emperor Antoninus, the passion of St. Pontian, 
 martyr, who was barbarously scourged for Christ by the command of the judge 
 Fabian, and then compelled to walk barefoot on burning coals. As he 
 was uninjured by the fire, he was put on the rack, was torn with iron hooks, 
 then thrown into a dungeon, where he was comforted by the visit of an angel. 
 He was afterwards exposed to the lions, had melted lead poured over him, and 
 finally died by the sword.
\switchcolumn*
\selectlanguage{latin}
Laudæ, in Insúbria, sancti Bassiáni, Epíscopi 
 et Confessóris; qui advérsus hæréticos, una cum sancto Ambrósio, strénue 
 decertávit.
\switchcolumn
\selectlanguage{english}
At Lodi in Lombardy, St. Bassian, bishop and confessor, who, in conjunction 
 with St. Ambrose, courageously combatted the heretics.
\switchcolumn*
\selectlanguage{latin}
Wigórniæ, in Anglia, sancti Wulstáni, Epíscopi 
 et Confessóris, méritis et miráculis conspícui; qui ab Innocéntio Papa 
 Tértio inter Sanctos relátus est.
\switchcolumn
\selectlanguage{english}
At Worcester, England, St. Wulfstan, bishop and confessor, conspicuous for 
 merits and miracles. He was ranked among the saints by Innocent III.
\switchcolumn*
\selectlanguage{latin}
\end{paracol}


% ---- martyrology/mart01/mart0120.htm
\needspace{10\baselineskip}
\begin{paracol}{2}
\selectlanguage{latin}
\begin{center}{\color{gregoriocolor} Tertiodécimo Kaléndas Februárii. 
 Luna\dots\ }\end{center}
\switchcolumn
\selectlanguage{english}
\begin{center}{\color{gregoriocolor} The 20\textsuperscript{th} Day of 
 January. The\dots\ Day of the Moon.}\end{center}
\end{paracol}

\noindent\begin{tabularx}{\linewidth}{*{19}{>{\centering\arraybackslash}X}}
 \textcolor{gregoriocolor}{a} & \textcolor{gregoriocolor}{b} & \textcolor{gregoriocolor}{c} & \textcolor{gregoriocolor}{d} & \textcolor{gregoriocolor}{e} & \textcolor{gregoriocolor}{f} & \textcolor{gregoriocolor}{g} & \textcolor{gregoriocolor}{h} & \textcolor{gregoriocolor}{i} & \textcolor{gregoriocolor}{k} & \textcolor{gregoriocolor}{l} & \textcolor{gregoriocolor}{m} & \textcolor{gregoriocolor}{n} & \textcolor{gregoriocolor}{p} & \textcolor{gregoriocolor}{q} & \textcolor{gregoriocolor}{r} & \textcolor{gregoriocolor}{s} & \textcolor{gregoriocolor}{t} & \textcolor{gregoriocolor}{u} \\
 21 & 22 & 23 & 24 & 25 & 26 & 27 & 28 & 29 & 30 & 1 & 2 & 3 & 4 & 5 & 6 & 7 & 8 & 9 \\
\end{tabularx}
\vspace{0.5\baselineskip}
\noindent\begin{tabularx}{\linewidth}{*{12}{>{\centering\arraybackslash}X}}
 \textcolor{gregoriocolor}{A} & \textcolor{gregoriocolor}{B} & \textcolor{gregoriocolor}{C} & \textcolor{gregoriocolor}{D} & \textcolor{gregoriocolor}{E} & F & \textcolor{gregoriocolor}{F} & \textcolor{gregoriocolor}{G} & \textcolor{gregoriocolor}{H} & \textcolor{gregoriocolor}{M} & \textcolor{gregoriocolor}{N} & \textcolor{gregoriocolor}{P} \\
 10 & 11 & 12 & 13 & 14 & 15 & 15 & 16 & 17 & 18 & 19 & 20 \\
\end{tabularx}

\begin{paracol}{2}
\selectlanguage{latin}
\lettrine[lines=2]{R}{omæ} natális sancti Fabiáni, Papæ et Mártyris; 
 qui, Décii témpore, martyrium passus est, atque in cœmetério Callísti 
 sepúltus.
\switchcolumn
\selectlanguage{english}
\lettrine[lines=2]{A}{t} Rome, the birthday of St. Fabian, pope, who suffered martyrdom in the 
 time of Decius, and was buried in the cemetery of Callistus.
\switchcolumn*
\selectlanguage{latin}
Item Romæ, ad Catacúmbas, sancti Sebastiáni 
 Mártyris, qui, Diocletiáno Imperatóre, cum habéret principátum primæ 
 cohórtis, jussus est, sub título christianitátis, ligári in médio campo, et 
 sagittári a milítibus, atque ad últimum fústibus cædi, donec defíceret.
\switchcolumn
\selectlanguage{english}
Also at Rome, in the catacombs, the martyr St. Sebastian. He was 
 commander of the first cohort under Emperor Diocletian, and for professing 
 Christianity he was bound to a tree in the centre of a vast field, shot with 
 arrows by the soldiers, and beaten with clubs until he expired.
\switchcolumn*
\selectlanguage{latin}
Nicǽæ, in Bithynia, sancti Neóphyti Mártyris, 
 qui, quintumdécimum annum ætátis agens, flagris cæsus, in fornácem immíssus, 
 feris objéctus, et, cum illæsus permanéret et Christi fidem constánter 
 profiterétur, gládio tandem occísus est.
\switchcolumn
\selectlanguage{english}
At Nicea in Bithynia, St. Neophytus, martyr, who in the fifteenth year of 
 his age, was scourged, cast into a furnace, and exposed to wild beasts. 
 As he remained uninjured, and constantly confessed the faith of Christ, he 
 was at last killed with the sword.
\switchcolumn*
\selectlanguage{latin}
Cæsénæ sancti Mauri Epíscopi, virtútibus et 
 miráculis clari.
\switchcolumn
\selectlanguage{english}
At Cesena, St. Maur, bishop, renowned for virtues and miracles.
\switchcolumn*
\selectlanguage{latin}
In Palæstína natális sancti Euthymii Abbátis, 
 qui zelo cathólicæ discíplinæ et virtúte miraculórum, témpore Marciáni 
 Imperatóris, in Ecclésia flóruit.
\switchcolumn
\selectlanguage{english}
In Palestine, in the time of Emperor Marcian, the birthday of St. Euthymius, 
 abbot, who flourished in the Church, full of zeal for Catholic discipline, 
 and gifted with miracles.
\switchcolumn*
\selectlanguage{latin}
\end{paracol}


% ---- martyrology/mart01/mart0121.htm
\needspace{10\baselineskip}
\begin{paracol}{2}
\selectlanguage{latin}
\begin{center}{\color{gregoriocolor} Duodécimo Kaléndas Februárii. 
 Luna\dots\ }\end{center}
\switchcolumn
\selectlanguage{english}
\begin{center}{\color{gregoriocolor} The 21\textsuperscript{st} Day of 
 January. The\dots\ Day of the Moon.}\end{center}
\end{paracol}

\noindent\begin{tabularx}{\linewidth}{*{19}{>{\centering\arraybackslash}X}}
 \textcolor{gregoriocolor}{a} & \textcolor{gregoriocolor}{b} & \textcolor{gregoriocolor}{c} & \textcolor{gregoriocolor}{d} & \textcolor{gregoriocolor}{e} & \textcolor{gregoriocolor}{f} & \textcolor{gregoriocolor}{g} & \textcolor{gregoriocolor}{h} & \textcolor{gregoriocolor}{i} & \textcolor{gregoriocolor}{k} & \textcolor{gregoriocolor}{l} & \textcolor{gregoriocolor}{m} & \textcolor{gregoriocolor}{n} & \textcolor{gregoriocolor}{p} & \textcolor{gregoriocolor}{q} & \textcolor{gregoriocolor}{r} & \textcolor{gregoriocolor}{s} & \textcolor{gregoriocolor}{t} & \textcolor{gregoriocolor}{u} \\
 22 & 23 & 24 & 25 & 26 & 27 & 28 & 29 & 30 & 1 & 2 & 3 & 4 & 5 & 6 & 7 & 8 & 9 & 10 \\
\end{tabularx}
\vspace{0.5\baselineskip}
\noindent\begin{tabularx}{\linewidth}{*{12}{>{\centering\arraybackslash}X}}
 \textcolor{gregoriocolor}{A} & \textcolor{gregoriocolor}{B} & \textcolor{gregoriocolor}{C} & \textcolor{gregoriocolor}{D} & \textcolor{gregoriocolor}{E} & F & \textcolor{gregoriocolor}{F} & \textcolor{gregoriocolor}{G} & \textcolor{gregoriocolor}{H} & \textcolor{gregoriocolor}{M} & \textcolor{gregoriocolor}{N} & \textcolor{gregoriocolor}{P} \\
 11 & 12 & 13 & 14 & 15 & 16 & 16 & 17 & 18 & 19 & 20 & 21 \\
\end{tabularx}

\begin{paracol}{2}
\selectlanguage{latin}
\lettrine[lines=2]{R}{omæ} pássio sanctæ Agnétis, Vírginis et 
 Mártyris; quæ, sub Præfécto Urbis Symphrónio, ígnibus injécta, sed iis per 
 oratiónem ejus exstínctis, gládio percússa est. De ea beátus 
 Hierónymus hæc scribit: « Omnium géntium lítteris atque linguis, præcípue in Ecclésiis, Agnétis vita laudáta est; quæ et ætátem vicit et tyránnum, et títulum castitátis martyrio consecrávit ».
\switchcolumn
\selectlanguage{english}
\lettrine[lines=2]{A}{t} Rome, the passion of St. Agnes, virgin, who under Symphronius, governor 
 of the city, was thrown into the fire, but after it was extinguished by her 
 prayers, she was slain with the sword. Of her, St. Jerome writes: 
 ``Agnes is praised in the writings and by the tongues of all nations, 
 especially in the churches. She overcame the weakness of her age, 
 conquered the cruelty of the tyrant, and consecrated her chastity by 
 martyrdom.''
\switchcolumn*
\selectlanguage{latin}
Athénis natális sancti Públii Epíscopi, qui Atheniénsium Ecclésiam, post 
 sanctum Dionysium Areopagítam, nobíliter rexit; et, præclárus 
 virtútibus ac doctrínæ laude præfúlgens, ob Christi martyrium glorióse 
 coronátur.
\switchcolumn
\selectlanguage{english}
At Athens, the birthday of St. Publius, bishop, who, as successor of St. 
 Denis the Areopagite, nobly governed the Church of Athens. No less 
 celebrated for the lustre of his virtues than for the brilliancy of his 
 learning, he was gloriously crowned for having borne testimony to Christ.
\switchcolumn*
\selectlanguage{latin}
Tarracóne, in Hispánia, sanctórum Mártyrum Fructuósi Epíscopi, Augúrii et 
 Eulógii Diaconórum. Hi, témpore Galliéni, primo in cárcerem trusi, 
 deínde flammis injécti, et, exústis vínculis, mánibus in modum crucis 
 expánsis orántes, martyrium complevérunt; in quorum die natáli sanctus 
 Augustínus sermónem ad pópulum hábuit.
\switchcolumn
\selectlanguage{english}
At Terragona in Spain, during the reign of Gallienus, the holy martyrs 
 Fructuosus, a bishop, Augurius and Eulogius, deacons. They were taken 
 from prison, cast into the fire, where, their bonds being burnt, they 
 extended their arms in the form of a cross, and thus in prayer they died. 
 On their anniversary, St. Augustine preached a sermon to his people.
\switchcolumn*
\selectlanguage{latin}
In monastério Einsidlénsi, apud Helvétios, sancti Meinrádi, Presbyteri et 
 Mónachi; qui eódem in loco, ubi póstea monastérium ipsum excrévit, eremíticæ 
 inténtus vitæ, a latrónibus interféctus est. Ipsíus vero beáti viri 
 corpus, olim in Augiénsi Germániæ monastério sepúltum, ad Einsidlénse 
 monastérium deínde relátum fuit.
\switchcolumn
\selectlanguage{english}
In the monastery of Einsiedeln in Switzerland, St. Meinrad, priest and monk, 
 who was slain by robbers after having lived as a hermit in this place where 
 the monastery was later built. The body of this holy man was first 
 buried in the monastery of Reichenau in Germany, and from there it was 
 transferred to the monastery of Einsiedeln.
\switchcolumn*
\selectlanguage{latin}
Trecis, in Gállia, sancti Pátrocli Mártyris, qui martyrii corónam sub 
 Aureliáno Imperatóre proméruit.
\switchcolumn
\selectlanguage{english}
At Troyes in France, St. Patroclus, martyr, who won the crown of martyrdom 
 under Emperor Aurelian.
\switchcolumn*
\selectlanguage{latin}
Papíæ sancti Epiphánii, Epíscopi et Confessóris.
\switchcolumn
\selectlanguage{english}
At Pavia, St. Epiphanius, bishop and confessor.
\switchcolumn*
\selectlanguage{latin}
\end{paracol}


% ---- martyrology/mart01/mart0122.htm
\needspace{10\baselineskip}
\begin{paracol}{2}
\selectlanguage{latin}
\begin{center}{\color{gregoriocolor} Undécimo Kaléndas Februárii. 
 Luna\dots\ }\end{center}
\switchcolumn
\selectlanguage{english}
\begin{center}{\color{gregoriocolor} The 22\textsuperscript{nd} Day of 
 January. The\dots\ Day of the Moon.}\end{center}
\end{paracol}

\noindent\begin{tabularx}{\linewidth}{*{19}{>{\centering\arraybackslash}X}}
 \textcolor{gregoriocolor}{a} & \textcolor{gregoriocolor}{b} & \textcolor{gregoriocolor}{c} & \textcolor{gregoriocolor}{d} & \textcolor{gregoriocolor}{e} & \textcolor{gregoriocolor}{f} & \textcolor{gregoriocolor}{g} & \textcolor{gregoriocolor}{h} & \textcolor{gregoriocolor}{i} & \textcolor{gregoriocolor}{k} & \textcolor{gregoriocolor}{l} & \textcolor{gregoriocolor}{m} & \textcolor{gregoriocolor}{n} & \textcolor{gregoriocolor}{p} & \textcolor{gregoriocolor}{q} & \textcolor{gregoriocolor}{r} & \textcolor{gregoriocolor}{s} & \textcolor{gregoriocolor}{t} & \textcolor{gregoriocolor}{u} \\
 23 & 24 & 25 & 26 & 27 & 28 & 29 & 30 & 1 & 2 & 3 & 4 & 5 & 6 & 7 & 8 & 9 & 10 & 11 \\
\end{tabularx}
\vspace{0.5\baselineskip}
\noindent\begin{tabularx}{\linewidth}{*{12}{>{\centering\arraybackslash}X}}
 \textcolor{gregoriocolor}{A} & \textcolor{gregoriocolor}{B} & \textcolor{gregoriocolor}{C} & \textcolor{gregoriocolor}{D} & \textcolor{gregoriocolor}{E} & F & \textcolor{gregoriocolor}{F} & \textcolor{gregoriocolor}{G} & \textcolor{gregoriocolor}{H} & \textcolor{gregoriocolor}{M} & \textcolor{gregoriocolor}{N} & \textcolor{gregoriocolor}{P} \\
 12 & 13 & 14 & 15 & 16 & 17 & 17 & 18 & 19 & 20 & 21 & 22 \\
\end{tabularx}

\begin{paracol}{2}
\selectlanguage{latin}
\lettrine[lines=2]{V}{aléntiæ,} in Hispánia Tarraconénsi, sancti 
 Vincéntii, Levítæ et Mártyris; qui, sub impiíssimo Præside Daciáno, cárceres, 
 famem, equúleum, distorsiónes membrórum, láminas candéntes, férream cratem 
 ignítam áliaque tormentórum génera perpéssus, ad martyrii præmium evolávit 
 in cælum; cujus passiónis nóbilem triúmphum Prudéntius luculénter vérsibus 
 exséquitur, et beátus Augustínus ac sanctus Leo Papa summis láudibus 
 comméndant.
\switchcolumn
\selectlanguage{english}
\lettrine[lines=2]{A}{t} Valencia in Spain, while the wicked Dacian was governor, St. Vincent, 
 deacon and martyr, who, after suffering imprisonment, hunger, the rack, and 
 the disjointing of his limbs, was burned with plates of heated metal and on 
 the gridiron, and tormented in other ways, then took his flight to heaven, 
 there to receive the reward of martyrdom. His noble triumph over his 
 sufferings has been skillfully set forth in verse by Prudentius, and also 
 was eulogized by St. Augustine and Pope St. Leo.
\switchcolumn*
\selectlanguage{latin}
Apud Bethsáloen, in Assyria, sancti Anastásii 
 Persæ Mónachi, qui, post plúrima torménta cárceris, vérberum et vinculórum, 
 quæ in Cæsaréa Palæstínæ perpéssus fúerat, a Persárum Rege Chósroa multis 
 pœnis afféctus, ad últimum decollátus est, cum prius septuagínta Sócios, qui 
 fúerant in fluénta demérsi, ad martyrium præmisísset. Ejus caput Romam, 
 ad Aquas Sálvias, delátum est, una cum veneránda ejus imágine, cujus aspéctu 
 fugári dæmones morbósque curári, Acta secúndi Concílii Nicǽni testántur.
\switchcolumn
\selectlanguage{english}
At Bethsaloen in Assyria, St. Anastasius, a Persian monk, who after 
 suffering much at Caesarea in Palestine from imprisonment, stripes, and 
 fetters, had to bear many afflictions from Chosroes, king of Persia, who 
 caused him to be beheaded. He had sent before him to martyrdom seventy 
 of his companions, who were drowned in a river. His head was brought 
 to Rome, at Aquae Salviae, together with his revered image, by the sight of 
 which demons are expelled, and diseases cured, as is attested by the Acts of 
 the second Council of Nicea.
\switchcolumn*
\selectlanguage{latin}
Ebredúni, in Gálliis, sanctórum Mártyrum Vincéntii, Oróntii et Victóris; qui 
 martyrio in Diocletiáni persecutióne coronáti sunt.
\switchcolumn
\selectlanguage{english}
At Embrun in France, the holy martyrs Vincent, Orontius, and Victor who were 
 crowned with martyrdom in the persecution of Diocletian.
\switchcolumn*
\selectlanguage{latin}
Nováriæ sancti Gaudéntii, Epíscopi et 
 Confessóris.
\switchcolumn
\selectlanguage{english}
At Novara, St. Gaudentius, bishop and confessor.
\switchcolumn*
\selectlanguage{latin}
Soræ sancti Domínici Abbátis, miráculis clari.
\switchcolumn
\selectlanguage{english}
At Sora, the abbot St. Dominic, renowned for miracles.
\switchcolumn*
\selectlanguage{latin}
\end{paracol}


% ---- martyrology/mart01/mart0123.htm
\needspace{10\baselineskip}
\begin{paracol}{2}
\selectlanguage{latin}
\begin{center}{\color{gregoriocolor} Décimo Kaléndas Februárii. 
 Luna\dots\ }\end{center}
\switchcolumn
\selectlanguage{english}
\begin{center}{\color{gregoriocolor} The 23\textsuperscript{rd} Day of 
 January. The\dots\ Day of the Moon.}\end{center}
\end{paracol}

\noindent\begin{tabularx}{\linewidth}{*{19}{>{\centering\arraybackslash}X}}
 \textcolor{gregoriocolor}{a} & \textcolor{gregoriocolor}{b} & \textcolor{gregoriocolor}{c} & \textcolor{gregoriocolor}{d} & \textcolor{gregoriocolor}{e} & \textcolor{gregoriocolor}{f} & \textcolor{gregoriocolor}{g} & \textcolor{gregoriocolor}{h} & \textcolor{gregoriocolor}{i} & \textcolor{gregoriocolor}{k} & \textcolor{gregoriocolor}{l} & \textcolor{gregoriocolor}{m} & \textcolor{gregoriocolor}{n} & \textcolor{gregoriocolor}{p} & \textcolor{gregoriocolor}{q} & \textcolor{gregoriocolor}{r} & \textcolor{gregoriocolor}{s} & \textcolor{gregoriocolor}{t} & \textcolor{gregoriocolor}{u} \\
 24 & 25 & 26 & 27 & 28 & 29 & 30 & 1 & 2 & 3 & 4 & 5 & 6 & 7 & 8 & 9 & 10 & 11 & 12 \\
\end{tabularx}
\vspace{0.5\baselineskip}
\noindent\begin{tabularx}{\linewidth}{*{12}{>{\centering\arraybackslash}X}}
 \textcolor{gregoriocolor}{A} & \textcolor{gregoriocolor}{B} & \textcolor{gregoriocolor}{C} & \textcolor{gregoriocolor}{D} & \textcolor{gregoriocolor}{E} & F & \textcolor{gregoriocolor}{F} & \textcolor{gregoriocolor}{G} & \textcolor{gregoriocolor}{H} & \textcolor{gregoriocolor}{M} & \textcolor{gregoriocolor}{N} & \textcolor{gregoriocolor}{P} \\
 13 & 14 & 15 & 16 & 17 & 18 & 18 & 19 & 20 & 21 & 22 & 23 \\
\end{tabularx}

\begin{paracol}{2}
\selectlanguage{latin}
\lettrine[lines=2]{S}{ancti} Raymúndi de Pénafort, ex Ordine Prædicatórum, 
 Confessóris; cujus dies natális octávo Idus Januárii recólitur.
\switchcolumn
\selectlanguage{english}
\lettrine[lines=2]{S}{t.} Raymond of Pennafort, of the Order of Preachers, whose birthday is the 
 sixth of this month.
\switchcolumn*
\selectlanguage{latin}
Romæ sanctæ Emerentiánæ, Vírginis et Mártyris; 
 quæ adhuc catechúmena, dum oráret ad sepúlcrum sanctæ Agnétis, cujus fúerat 
 collactánea, a Gentílibus lapidáta est.
\switchcolumn
\selectlanguage{english}
At Rome, the holy virgin and martyr, St. Emerentiana. Being yet a 
 catechumen, she was stoned to death by the heathens while praying at the 
 tomb of St. Agnes, her foster sister.
\switchcolumn*
\selectlanguage{latin}
Philíppis, in Macedónia, sancti Pármenæ, qui 
 fuit unus de septem primis Diáconis. Hic, tráditus grátiæ Dei, injúnctum sibi a frátribus offícium prædicatiónis plena fide consúmmans, 
 martyrii glóriam, sub Trajáno, est adéptus.
\switchcolumn
\selectlanguage{english}
At Philippi in Macedonia, St. Parmenas, one of the first seven deacons, who 
 by the grace of God faithfully discharged the office of preaching committed 
 to him, and obtained the glory of martyrdom in the time of Trajan.
\switchcolumn*
\selectlanguage{latin}
Ancyræ, in Galátia, sancti Cleméntis Epíscopi, 
 qui, sæpius cruciátus, tandem, sub Diocletiáno Imperatóre, martyrium 
 consummávit.
\switchcolumn
\selectlanguage{english}
At Ancyra in Galatia, St. Clement, bishop. After enduring frequent 
 torments, he finally completed his martyrdom under Diocletian.
\switchcolumn*
\selectlanguage{latin}
Ibídem sancti Agathángeli, qui eódem die, sub 
 Lúcio Præside, passus est.
\switchcolumn
\selectlanguage{english}
In the same place, and on the same day, St. Agathangelus who suffered under 
 the governor Lucius.
\switchcolumn*
\selectlanguage{latin}
Cæsaréæ, in Mauritánia, sanctórum Mártyrum 
 Severiáni et Aquilæ uxóris, ígnibus combustórum.
\switchcolumn
\selectlanguage{english}
At Caesarea in Morocco, the holy martyrs Severian and his wife Aquila, who 
 were consumed by fire.
\switchcolumn*
\selectlanguage{latin}
Apud Antínoum, Ægypti urbem, sancti Asclæ 
 Mártyris, qui, post divérsa torménta, pretiósam Deo ánimam, in flumen 
 præcipitátus, réddidit.
\switchcolumn
\selectlanguage{english}
At Antinoum, a city of Egypt, St. Ascla, martyr, who, after various 
 torments, was thrown into a river and gave up his precious soul unto God.
\switchcolumn*
\selectlanguage{latin}
Alexandríæ sancti Joánnis Eleemosynárii, 
 ejúsdem urbis Epíscopi, misericórdia in páuperes celebérrimi.
\switchcolumn
\selectlanguage{english}
At Alexandria, St. John the Almoner, bishop of that city, celebrated for his 
 charity towards the poor.
\switchcolumn*
\selectlanguage{latin}
Toléti, in Hispánia, sancti Ildefónsi Epíscopi, qui, ob singulárem vitæ 
 integritátem, susceptámque fídei defensiónem advérsus hæréticos, sanctíssimæ 
 Dei Genitrícis virginitátem impugnántes, ab eádem Vírgine María donátus est 
 candidíssima veste, ac demum, sanctitáte célebris, in cælum vocátus.
\switchcolumn
\selectlanguage{english}
At Toledo, St. Ildefonse, bishop, renowned for sanctity. On account of 
 his great purity of life, and his defence of the virginity of the Mother of 
 God, against the heretics who denied it, he received from her a brilliant 
 white vestment, and was called to heaven.
\switchcolumn*
\selectlanguage{latin}
In Província Valériæ sancti Martyrii Mónachi, 
 cujus méminit beátus Gregórius Papa.
\switchcolumn
\selectlanguage{english}
In the province of Valeria, St. Martyrius, monk, mentioned by Pope St. 
 Gregory.
\switchcolumn*
\selectlanguage{latin}
\end{paracol}


% ---- martyrology/mart01/mart0124.htm
\needspace{10\baselineskip}
\begin{paracol}{2}
\selectlanguage{latin}
\begin{center}{\color{gregoriocolor} Nono Kaléndas Februárii. 
 Luna\dots\ }\end{center}
\switchcolumn
\selectlanguage{english}
\begin{center}{\color{gregoriocolor} The 24\textsuperscript{th} Day of 
 January. The\dots\ Day of the Moon.}\end{center}
\end{paracol}

\noindent\begin{tabularx}{\linewidth}{*{19}{>{\centering\arraybackslash}X}}
 \textcolor{gregoriocolor}{a} & \textcolor{gregoriocolor}{b} & \textcolor{gregoriocolor}{c} & \textcolor{gregoriocolor}{d} & \textcolor{gregoriocolor}{e} & \textcolor{gregoriocolor}{f} & \textcolor{gregoriocolor}{g} & \textcolor{gregoriocolor}{h} & \textcolor{gregoriocolor}{i} & \textcolor{gregoriocolor}{k} & \textcolor{gregoriocolor}{l} & \textcolor{gregoriocolor}{m} & \textcolor{gregoriocolor}{n} & \textcolor{gregoriocolor}{p} & \textcolor{gregoriocolor}{q} & \textcolor{gregoriocolor}{r} & \textcolor{gregoriocolor}{s} & \textcolor{gregoriocolor}{t} & \textcolor{gregoriocolor}{u} \\
 25 & 26 & 27 & 28 & 29 & 30 & 1 & 2 & 3 & 4 & 5 & 6 & 7 & 8 & 9 & 10 & 11 & 12 & 13 \\
\end{tabularx}
\vspace{0.5\baselineskip}
\noindent\begin{tabularx}{\linewidth}{*{12}{>{\centering\arraybackslash}X}}
 \textcolor{gregoriocolor}{A} & \textcolor{gregoriocolor}{B} & \textcolor{gregoriocolor}{C} & \textcolor{gregoriocolor}{D} & \textcolor{gregoriocolor}{E} & F & \textcolor{gregoriocolor}{F} & \textcolor{gregoriocolor}{G} & \textcolor{gregoriocolor}{H} & \textcolor{gregoriocolor}{M} & \textcolor{gregoriocolor}{N} & \textcolor{gregoriocolor}{P} \\
 14 & 15 & 16 & 17 & 18 & 19 & 19 & 20 & 21 & 22 & 23 & 24 \\
\end{tabularx}

\begin{paracol}{2}
\selectlanguage{latin}
\lettrine[lines=2]{A}{pud} Ephesum sancti Timóthei, qui fuit 
 discípulus beáti Pauli Apóstoli; atque, ab eódem Ephesi ordinátus Epíscopus, 
 ibi, post multos pro Christo agónes, cum Diánæ immolántes argúeret, 
 lapídibus óbrutus est, ac paulo post obdormívit in Dómino.
\switchcolumn
\selectlanguage{english}
\lettrine[lines=2]{A}{t} Ephesus, St. Timothy, disciple of the apostle St. Paul, who ordained him 
 bishop of that city. After many labours for Christ, he was stoned for 
 rebuking those who offered sacrifices to Diana, and shortly after went 
 peacefully to his rest in the Lord.
\switchcolumn*
\selectlanguage{latin}
Antiochíæ sancti Bábilæ Epíscopi, qui, in 
 persecutióne Décii, póstea quam frequénter passiónibus suis ac cruciátibus 
 glorificáverat Deum, gloriósæ vitæ finem sortítus est in vínculis férreis, 
 cum quibus et suum corpus sepelíri mandávit. Referúntur étiam passi 
 cum eo tres púeri, scílicet Urbánus, Prilidiánus et Epolónius, quos ille in 
 Christi fide instrúxerat.
\switchcolumn
\selectlanguage{english}
At Antioch, in the persecution of Decius, Bishop St. Babylas, who frequently 
 glorified God by his sufferings and torments, ended his life in chains, with 
 which he ordered his body to be buried. Three boys, whom he had 
 instructed in the faith of Christ, Urbanus, Prilidian, and Epolonius, are 
 said to have suffered with him.
\switchcolumn*
\selectlanguage{latin}
Fulgínei, in Umbria, sancti Feliciáni, qui, a sancto Victóre Papa Primo 
 Epíscopus ejúsdem civitátis ordinátus, illic, post multos labóres, in 
 última senectúte, sub Décio Imperatóre, martyrio coronátus est.
\switchcolumn
\selectlanguage{english}
At Foligno in Umbria, St. Felician, consecrated bishop of that city by Pope 
 St. Victor I. After many labours, in extreme old age, he was crowned 
 with martyrdom in the time of Decius.
\switchcolumn*
\selectlanguage{latin}
Neocæsaréæ, in Mauritánia, sanctórum Mártyrum 
 Mardónii, Musónii, Eugénii et Metélli; qui omnes igni tráditi sunt, et eórum 
 relíquiæ in flumen dispérsæ.
\switchcolumn
\selectlanguage{english}
At Neocaesarea, the holy martyrs Mardonius, Musonius, Eugenius, and Metellus, 
 who were all burned to death, and their remains thrown into the river.
\switchcolumn*
\selectlanguage{latin}
Item sanctórum Mártyrum Thyrsi et Projécti.
\switchcolumn
\selectlanguage{english}
Also, the holy martyrs Thyrsus and Projectus.
\switchcolumn*
\selectlanguage{latin}
Cínguli, in Picéno, sancti Exsuperántii 
 Confessóris, ejúsdem civitátis Epíscopi, ob miraculórum famam illústris.
\switchcolumn
\selectlanguage{english}
At Cingoli in Piceno, St. Exuperantius, confessor and bishop of that city, 
 who attained great fame by his miracles.
\switchcolumn*
\selectlanguage{latin}
Bonóniæ sancti Zamæ, qui, a sancto Dionysio, 
 Románo Pontífice, primus ejúsdem civitátis Epíscopus ordinátus, illic 
 Christiánam fidem mirífice propagávit.
\switchcolumn
\selectlanguage{english}
At Bologna, St. Zamas, the first bishop of that city, who was consecrated by 
 Pope St. Denis, and there did wonders in spreading the Christian faith.
\switchcolumn*
\selectlanguage{latin}
Item beáti Suráni Abbátis, qui, témpore Longobardórum, sanctitáte flóruit.
\switchcolumn
\selectlanguage{english}
Also, blessed Suranus, abbot, who lived in the time of the Lombards.
\switchcolumn*
\selectlanguage{latin}
\end{paracol}


% ---- martyrology/mart01/mart0125.htm
\needspace{10\baselineskip}
\begin{paracol}{2}
\selectlanguage{latin}
\begin{center}{\color{gregoriocolor} Octávo Kaléndas Februárii. 
 Luna\dots\ }\end{center}
\switchcolumn
\selectlanguage{english}
\begin{center}{\color{gregoriocolor} The 25\textsuperscript{th} Day of 
 January. The\dots\ Day of the Moon.}\end{center}
\end{paracol}

\noindent\begin{tabularx}{\linewidth}{*{19}{>{\centering\arraybackslash}X}}
 \textcolor{gregoriocolor}{a} & \textcolor{gregoriocolor}{b} & \textcolor{gregoriocolor}{c} & \textcolor{gregoriocolor}{d} & \textcolor{gregoriocolor}{e} & \textcolor{gregoriocolor}{f} & \textcolor{gregoriocolor}{g} & \textcolor{gregoriocolor}{h} & \textcolor{gregoriocolor}{i} & \textcolor{gregoriocolor}{k} & \textcolor{gregoriocolor}{l} & \textcolor{gregoriocolor}{m} & \textcolor{gregoriocolor}{n} & \textcolor{gregoriocolor}{p} & \textcolor{gregoriocolor}{q} & \textcolor{gregoriocolor}{r} & \textcolor{gregoriocolor}{s} & \textcolor{gregoriocolor}{t} & \textcolor{gregoriocolor}{u} \\
 26 & 27 & 28 & 29 & 30 & 1 & 2 & 3 & 4 & 5 & 6 & 7 & 8 & 9 & 10 & 11 & 12 & 13 & 14 \\
\end{tabularx}
\vspace{0.5\baselineskip}
\noindent\begin{tabularx}{\linewidth}{*{12}{>{\centering\arraybackslash}X}}
 \textcolor{gregoriocolor}{A} & \textcolor{gregoriocolor}{B} & \textcolor{gregoriocolor}{C} & \textcolor{gregoriocolor}{D} & \textcolor{gregoriocolor}{E} & F & \textcolor{gregoriocolor}{F} & \textcolor{gregoriocolor}{G} & \textcolor{gregoriocolor}{H} & \textcolor{gregoriocolor}{M} & \textcolor{gregoriocolor}{N} & \textcolor{gregoriocolor}{P} \\
 15 & 16 & 17 & 18 & 19 & 20 & 20 & 21 & 22 & 23 & 24 & 25 \\
\end{tabularx}

\begin{paracol}{2}
\selectlanguage{latin}
\lettrine[lines=2]{C}{onvérsio} sancti Pauli Apóstoli, quæ fuit anno 
 secúndo ab Ascensióne Domini.
\switchcolumn
\selectlanguage{english}
\lettrine[lines=2]{T}{he} conversion of St. Paul the Apostle, which occurred in the second year 
 after the Ascension of our Lord.
\switchcolumn*
\selectlanguage{latin}
Apud Damáscum natális sancti Ananíæ, qui fuit 
 discípulus Dómini, et eúndem Paulum Apóstolum baptizávit. Ipse autem, 
 cum Damásci, et Eleutherópoli, alibíque Evangélium prædicásset, tandem, sub 
 Licínio Júdice, nervis cæsus et laniátus, ac lapídibus oppréssus, martyrium 
 consummávit.
\switchcolumn
\selectlanguage{english}
At Damascus, the birthday of St. Ananias, who was a disciple of our Lord, 
 and baptized the apostle Paul. After he had preached the Gospel at 
 Damascus, Eleutheropolis, and elsewhere, he was scourged under the judge 
 Licinius, had his flesh torn, and lastly being overwhelmed with stones, 
 ended his martyrdom.
\switchcolumn*
\selectlanguage{latin}
Arvérnis, in Gállia, sanctórum Præjécti 
 Epíscopi, et Amaríni, Abbátis Cloroangiénsis; qui ambo a procéribus ejúsdem 
 urbis passi sunt.
\switchcolumn
\selectlanguage{english}
In the Auvergne in France, the Saints Praejectus, bishop, and Amarinus, 
 abbot of Doroang, who were murdered by the leading men of that city.
\switchcolumn*
\selectlanguage{latin}
Antiochíæ sanctórum Mártyrum Juventíni et 
 Máximi, qui, sub Juliáno Apóstata, martyrio coronáti sunt; in quorum die 
 natáli sanctus Joánnes Chrysóstomus sermónem ad pópulum hábuit.
\switchcolumn
\selectlanguage{english}
At Antioch, in the time of Julian the Apostate, the holy martyrs Juvenius 
 and Maximus, who were crowned with martyrdom. On their birthday, St. 
 John Chrysostom preached a sermon to his people.
\switchcolumn*
\selectlanguage{latin}
Item sanctórum Mártyrum Donáti, Sabíni et Agapis.
\switchcolumn
\selectlanguage{english}
Also, the holy martyrs Donatus, Sabinus, and Agape.
\switchcolumn*
\selectlanguage{latin}
Tomis, in Scythia, sancti Bretanniónis Epíscopi, qui mira sanctitáte et 
 cathólicæ fídei zelo, sub Ariáno Imperatóre 
 Valénte, cui fórtiter réstitit, in Ecclésia flóruit.
\switchcolumn
\selectlanguage{english}
At Tomis in Scythia, St. Bretannio, bishop, who worked in the Church shewing 
 great sanctity and zeal for the Catholic faith, and was at the same time 
 bravely opposed to the Arian emperor Valens.
\switchcolumn*
\selectlanguage{latin}
Marciánis, in Gállia, sancti Poppónis, Presbyteri et Abbátis, miráculis 
 clari.
\switchcolumn
\selectlanguage{english}
At Marchiennes in France, St. Poppo, priest and abbot, renowned for his 
 miracles.
\switchcolumn*
\selectlanguage{latin}
\end{paracol}


% ---- martyrology/mart01/mart0126.htm
\needspace{10\baselineskip}
\begin{paracol}{2}
\selectlanguage{latin}
\begin{center}{\color{gregoriocolor} Séptimo Kaléndas Februárii. 
 Luna\dots\ }\end{center}
\switchcolumn
\selectlanguage{english}
\begin{center}{\color{gregoriocolor} The 26\textsuperscript{th} Day of 
 January. The\dots\ Day of the Moon.}\end{center}
\end{paracol}

\noindent\begin{tabularx}{\linewidth}{*{19}{>{\centering\arraybackslash}X}}
 \textcolor{gregoriocolor}{a} & \textcolor{gregoriocolor}{b} & \textcolor{gregoriocolor}{c} & \textcolor{gregoriocolor}{d} & \textcolor{gregoriocolor}{e} & \textcolor{gregoriocolor}{f} & \textcolor{gregoriocolor}{g} & \textcolor{gregoriocolor}{h} & \textcolor{gregoriocolor}{i} & \textcolor{gregoriocolor}{k} & \textcolor{gregoriocolor}{l} & \textcolor{gregoriocolor}{m} & \textcolor{gregoriocolor}{n} & \textcolor{gregoriocolor}{p} & \textcolor{gregoriocolor}{q} & \textcolor{gregoriocolor}{r} & \textcolor{gregoriocolor}{s} & \textcolor{gregoriocolor}{t} & \textcolor{gregoriocolor}{u} \\
 27 & 28 & 29 & 30 & 1 & 2 & 3 & 4 & 5 & 6 & 7 & 8 & 9 & 10 & 11 & 12 & 13 & 14 & 15 \\
\end{tabularx}
\vspace{0.5\baselineskip}
\noindent\begin{tabularx}{\linewidth}{*{12}{>{\centering\arraybackslash}X}}
 \textcolor{gregoriocolor}{A} & \textcolor{gregoriocolor}{B} & \textcolor{gregoriocolor}{C} & \textcolor{gregoriocolor}{D} & \textcolor{gregoriocolor}{E} & F & \textcolor{gregoriocolor}{F} & \textcolor{gregoriocolor}{G} & \textcolor{gregoriocolor}{H} & \textcolor{gregoriocolor}{M} & \textcolor{gregoriocolor}{N} & \textcolor{gregoriocolor}{P} \\
 16 & 17 & 18 & 19 & 20 & 21 & 21 & 22 & 23 & 24 & 25 & 26 \\
\end{tabularx}

\begin{paracol}{2}
\selectlanguage{latin}
\lettrine[lines=2]{S}{ancti} Polycárpi, Epíscopi Smyrnénsis et Mártyris; qui martyrii corónam 
 séptimo Kaléndas Mártii consecútus est.
\switchcolumn
\selectlanguage{english}
\lettrine[lines=2]{S}{t.} Polycarp, Bishop of Smyrna and martyr, who gained the crown of martyrdom 
 on the 23rd of February.
\switchcolumn*
\selectlanguage{latin}
Hippóne Régio, in Africa, sanctórum Theógenis 
 Epíscopi, et aliórum trigínta sex; qui, in persecutióne Valeriáni, 
 contemnéntes temporálem mortem, corónam ætérnæ vitæ adépti sunt.
\switchcolumn
\selectlanguage{english}
At Hippo in Africa, the holy bishop Theogenes and thirty-six others, who, 
 despising temporal death, obtained the crown of eternal life in the 
 persecution of Valerian.
\switchcolumn*
\selectlanguage{latin}
Apud Béthlehem Judæ dormítio sanctæ Paulæ Víduæ, 
 quæ, cum esset e nobilíssimo Senatórum génere, cum beáta Vírgine Christi 
 Eustóchio, fília sua, renúntians sæculo, facultátes suas paupéribus 
 distríbuit, et ad Præsépe Dómini se recépit; ibíque, multis virtútibus 
 prǽdita et longo coronáta martyrio, ad cæléstia regna transívit. 
 Ipsíus autem vitam, virtútibus admirándum, sanctus Hierónymus scripsit.
\switchcolumn
\selectlanguage{english}
At Bethlehem of Judea, the death of St. Paula, widow, mother of St. 
 Eustochium, a virgin of Christ, who abandoned her worldly prospects, though 
 she was descended from a noble line of senators, distributed her goods to 
 the poor, and retired to our Lord's manger, where, endowed with many 
 virtues, and crowned with a long martyrdom, she departed for the kingdom of 
 heaven. Her admirable life was written by St. Jerome.
\switchcolumn*
\selectlanguage{latin}
\end{paracol}


% ---- martyrology/mart01/mart0127.htm
\needspace{10\baselineskip}
\begin{paracol}{2}
\selectlanguage{latin}
\begin{center}{\color{gregoriocolor} Sexto Kaléndas Februárii. 
 Luna\dots\ }\end{center}
\switchcolumn
\selectlanguage{english}
\begin{center}{\color{gregoriocolor} The 27\textsuperscript{th} Day of 
 January. The\dots\ Day of the Moon.}\end{center}
\end{paracol}

\noindent\begin{tabularx}{\linewidth}{*{19}{>{\centering\arraybackslash}X}}
 \textcolor{gregoriocolor}{a} & \textcolor{gregoriocolor}{b} & \textcolor{gregoriocolor}{c} & \textcolor{gregoriocolor}{d} & \textcolor{gregoriocolor}{e} & \textcolor{gregoriocolor}{f} & \textcolor{gregoriocolor}{g} & \textcolor{gregoriocolor}{h} & \textcolor{gregoriocolor}{i} & \textcolor{gregoriocolor}{k} & \textcolor{gregoriocolor}{l} & \textcolor{gregoriocolor}{m} & \textcolor{gregoriocolor}{n} & \textcolor{gregoriocolor}{p} & \textcolor{gregoriocolor}{q} & \textcolor{gregoriocolor}{r} & \textcolor{gregoriocolor}{s} & \textcolor{gregoriocolor}{t} & \textcolor{gregoriocolor}{u} \\
 28 & 29 & 30 & 1 & 2 & 3 & 4 & 5 & 6 & 7 & 8 & 9 & 10 & 11 & 12 & 13 & 14 & 15 & 16 \\
\end{tabularx}
\vspace{0.5\baselineskip}
\noindent\begin{tabularx}{\linewidth}{*{12}{>{\centering\arraybackslash}X}}
 \textcolor{gregoriocolor}{A} & \textcolor{gregoriocolor}{B} & \textcolor{gregoriocolor}{C} & \textcolor{gregoriocolor}{D} & \textcolor{gregoriocolor}{E} & F & \textcolor{gregoriocolor}{F} & \textcolor{gregoriocolor}{G} & \textcolor{gregoriocolor}{H} & \textcolor{gregoriocolor}{M} & \textcolor{gregoriocolor}{N} & \textcolor{gregoriocolor}{P} \\
 17 & 18 & 19 & 20 & 21 & 22 & 22 & 23 & 24 & 25 & 26 & 27 \\
\end{tabularx}

\begin{paracol}{2}
\selectlanguage{latin}
\lettrine[lines=2]{S}{ancti} Joánnis Chrysóstomi, Epíscopi Constantinopolitáni, Confessóris et 
 Ecclésiæ Doctóris, cæléstis Oratórum sacrórum 
 Patróni; qui décimo octávo Kaléndas Octóbris obdormívit in Dómino. 
 Ejus sacrum corpus, sub Theodósio junióre, hac die Constantinópolim, inde 
 póstea Romam translátum fuit, et in Basílica Príncipis Apostolórum cónditum.
\switchcolumn
\selectlanguage{english}
\lettrine[lines=2]{S}{t.} John Chrysostom, Bishop of Constantinople, confessor and doctor of the 
 Church, and the heavenly patron of preachers, who fell asleep in the Lord on 
 the 14th of September. His holy body was brought to Constantinople on 
 this day in the reign of Theodosius the younger; it was afterwards taken to 
 Rome and placed in the basilica of the Prince of the Apostles.
\switchcolumn*
\selectlanguage{latin}
Bríxiæ natális sanctæ Angelæ Meríci Vírginis, 
 ex tértio Ordine sancti Francísci, quæ Societátem Vírginum sanctæ Ursulæ 
 instítuit, quarum præcípuum munus esset dirígere adolescéntulas in vias 
 Dómini. Ejus tamen festívitas 
 Kaléndis Júnii celebrátur.
\switchcolumn
\selectlanguage{english}
At Brescia, the birthday of St. Angela Merici, virgin, who belonged to the 
 Third Order of St. Francis, and who founded the Order of the Nuns of St. 
 Ursula, whose principal aim is to direct young girls in the ways of the 
 Lord. Her feast is celebrated on the 1st 
 of June.
\switchcolumn*
\selectlanguage{latin}
Apud Cenómanos, in Gállia, deposítio sancti 
 Juliáni, ejúsdem urbis primi Epíscopi, quem sanctus Petrus illuc ad 
 prædicándum Evangélium misit.
\switchcolumn
\selectlanguage{english}
At Le Mans in France, the death of St. Julian, the first bishop of that 
 city, who was sent there by St. Peter to preach the Gospel.
\switchcolumn*
\selectlanguage{latin}
Soræ sancti Juliáni Mártyris, qui, in 
 persecutióne Antoníni, sub Flaviáno Præside, comprehénsus est, et, cum 
 idolórum templum, dum ipse torquerétur, corruísset, martyrii corónam, 
 truncáto cápite, accépit.
\switchcolumn
\selectlanguage{english}
At Sora, St. Julian, martyr, who, being arrested in the persecution of 
 Antoninus, was beheaded because a pagan temple had fallen to the ground 
 while he was being tortured.
\switchcolumn*
\selectlanguage{latin}
In Africa sancti Avíti Mártyris.
\switchcolumn
\selectlanguage{english}
In Africa, St. Avitus, martyr.
\switchcolumn*
\selectlanguage{latin}
Ibídem sanctórum Mártyrum Dátii, Reátri et Sociórum, qui in persecutióne 
 Wandálica passi sunt.
\switchcolumn
\selectlanguage{english}
In the same country, the holy martyrs Datius, Reatrus, and their companions, 
 who suffered in the persecution of the Vandals.
\switchcolumn*
\selectlanguage{latin}
Item sanctórum Datívi, Juliáni, Vincéntii atque aliórum vigínti septem 
 Mártyrum.
\switchcolumn
\selectlanguage{english}
Also, the holy martyrs Dativus, Julian, Vincent, and twenty-seven others.
\switchcolumn*
\selectlanguage{latin}
Romæ sancti Vitaliáni Papæ.
\switchcolumn
\selectlanguage{english}
At Rome, St. Vitalian, pope.
\switchcolumn*
\selectlanguage{latin}
In monastério Bobacénsi, in Gállia, sancti Mauri Abbátis.
\switchcolumn
\selectlanguage{english}
In the monastery of Bobbio in France, St. Maur, abbot.
\switchcolumn*
\selectlanguage{latin}
\end{paracol}


% ---- martyrology/mart01/mart0128.htm
\needspace{10\baselineskip}
\begin{paracol}{2}
\selectlanguage{latin}
\begin{center}{\color{gregoriocolor} Quinto Kaléndas Februárii. 
 Luna\dots\ }\end{center}
\switchcolumn
\selectlanguage{english}
\begin{center}{\color{gregoriocolor} The 28\textsuperscript{th} Day of 
 January. The\dots\ Day of the Moon.}\end{center}
\end{paracol}

\noindent\begin{tabularx}{\linewidth}{*{19}{>{\centering\arraybackslash}X}}
 \textcolor{gregoriocolor}{a} & \textcolor{gregoriocolor}{b} & \textcolor{gregoriocolor}{c} & \textcolor{gregoriocolor}{d} & \textcolor{gregoriocolor}{e} & \textcolor{gregoriocolor}{f} & \textcolor{gregoriocolor}{g} & \textcolor{gregoriocolor}{h} & \textcolor{gregoriocolor}{i} & \textcolor{gregoriocolor}{k} & \textcolor{gregoriocolor}{l} & \textcolor{gregoriocolor}{m} & \textcolor{gregoriocolor}{n} & \textcolor{gregoriocolor}{p} & \textcolor{gregoriocolor}{q} & \textcolor{gregoriocolor}{r} & \textcolor{gregoriocolor}{s} & \textcolor{gregoriocolor}{t} & \textcolor{gregoriocolor}{u} \\
 29 & 30 & 1 & 2 & 3 & 4 & 5 & 6 & 7 & 8 & 9 & 10 & 11 & 12 & 13 & 14 & 15 & 16 & 17 \\
\end{tabularx}
\vspace{0.5\baselineskip}
\noindent\begin{tabularx}{\linewidth}{*{12}{>{\centering\arraybackslash}X}}
 \textcolor{gregoriocolor}{A} & \textcolor{gregoriocolor}{B} & \textcolor{gregoriocolor}{C} & \textcolor{gregoriocolor}{D} & \textcolor{gregoriocolor}{E} & F & \textcolor{gregoriocolor}{F} & \textcolor{gregoriocolor}{G} & \textcolor{gregoriocolor}{H} & \textcolor{gregoriocolor}{M} & \textcolor{gregoriocolor}{N} & \textcolor{gregoriocolor}{P} \\
 18 & 19 & 20 & 21 & 22 & 23 & 23 & 24 & 25 & 26 & 27 & 28 \\
\end{tabularx}

\begin{paracol}{2}
\selectlanguage{latin}
\lettrine[lines=2]{S}{ancti} Petri Nolásci Confessóris, qui Ordinis beátæ 
 Maríæ de Mercéde redemptiónis captivórum éxstitit Fundátor, et octávo Kaléndas Januárii obdormívit in Dómino.
\switchcolumn
\selectlanguage{english}
\lettrine[lines=2]{S}{t.} Peter Nolasco, confessor, who founded the Order of Our Lady of Ransom 
 for the redemption of captives, and who fell asleep in the Lord on the 25th 
 of December.
\switchcolumn*
\selectlanguage{latin}
Romæ sanctæ Agnétis, Vírginis et Mártyris, 
 secúndo.
\switchcolumn
\selectlanguage{english}
At Rome, the second feast of St. Agnes, virgin and martyr.
\switchcolumn*
\selectlanguage{latin}
Alexandríæ natális sancti Cyrílli, ejúsdem 
 urbis Epíscopi, Confessóris et Ecclésiæ Doctóris; qui, cathólicæ fídei 
 præclaríssimus propugnátor, doctrína et sanctitáte illústris quiévit in 
 pace. Ejus tamen festívitas quinto Idus Februárii celebrátur.
\switchcolumn
\selectlanguage{english}
At Alexandria, the birthday of St. Cyril, bishop of that city, a most 
 celebrated defender of the Catholic faith, who died in peace, with a great 
 reputation for learning and sanctity. His feast, however, is kept on 
 the ninth of February.
\switchcolumn*
\selectlanguage{latin}
Romæ sancti Flaviáni Mártyris, qui sub 
 Diocletiáno passus est.
\switchcolumn
\selectlanguage{english}
At Rome, St. Flavian, martyr, who suffered under Diocletian.
\switchcolumn*
\selectlanguage{latin}
Alexandríæ pássio plurimórum sanctórum Mártyrum, 
 qui, hac ipsa die, a factióne Syriáni, Ducis Ariáni, dum in Ecclésia synáxim 
 ágerent, divérso mortis génere sunt interémpti.
\switchcolumn
\selectlanguage{english}
At Alexandria, the commemoration of many holy martyrs, who, while they were 
 at Mass in the church on this day, were put to death in different ways by 
 the followers of Syrianus, an Arian general.
\switchcolumn*
\selectlanguage{latin}
Apollóniæ sanctórum Mártyrum, Léucii, Thyrsi et 
 Calliníci; qui, témpore Décii Imperatóris, váriis tormentórum genéribus 
 cruciáti, ac primus et últimus abscissióne cápitis, médius cælésti voce 
 evocátus spíritum reddens, martyrium consummárunt.
\switchcolumn
\selectlanguage{english}
At Appollonia, the holy martyrs Thrysus, Leucius, and Callinicus, who were 
 made to undergo various torments in the time of Emperor Decius. 
 Thyrsus and Callinicus were beheaded; Leucius, called by a heavenly voice, 
 yielded his soul unto God.
\switchcolumn*
\selectlanguage{latin}
In Thebáide sanctórum Mártyrum Leónidæ et 
 Sociórum, qui, témpore Diocletiáni, palmam martyrii sunt assecúti.
\switchcolumn
\selectlanguage{english}
In Thebais, the holy martyrs Leonides and his companions, who obtained the 
 palm of martyrdom in the time of Diocletian.
\switchcolumn*
\selectlanguage{latin}
Cæsaraugústæ, in Hispánia, sancti Valérii 
 Epíscopi.
\switchcolumn
\selectlanguage{english}
At Saragossa in Spain, St. Valerius, bishop.
\switchcolumn*
\selectlanguage{latin}
Conchæ, in Hispánia, natális sancti Juliáni 
 Epíscopi, qui, érogans in páuperes bona Ecclésiæ, ópera mánuum sibi more 
 Apostólico victum quærens, clarus miráculis quiévit in pace.
\switchcolumn
\selectlanguage{english}
At Cuenca in Spain, the birthday of St. Julian, bishop, who, after bestowing 
 the goods of the Church on the poor, like the apostles, supported himself by the work of 
 his hands, and went to his God famous for his miracles.
\switchcolumn*
\selectlanguage{latin}
In monastério Reomaénsi, in Gállia, deposítio sancti Joánnis Presbyteri, 
 viri Deo devóti.
\switchcolumn
\selectlanguage{english}
In the monastery of Rheims in France, the death of the holy priest John, a 
 devout man of God.
\switchcolumn*
\selectlanguage{latin}
In Palæstína sancti Jacóbi Eremítæ, qui, post 
 lapsum, diu, pæniténtiæ causa, in sepúlcro látuit, et clarus miráculis 
 migrávit ad Dóminum.
\switchcolumn
\selectlanguage{english}
In Palestine, St. James, hermit, who hid himself a long time in a sepulchre 
 in order to do penance for a fault he had committed, and, being celebrated 
 for miracles, departed for heaven.
\switchcolumn*
\selectlanguage{latin}
\end{paracol}


% ---- martyrology/mart01/mart0129.htm
\needspace{10\baselineskip}
\begin{paracol}{2}
\selectlanguage{latin}
\begin{center}{\color{gregoriocolor} Quarto Kaléndas Februárii. 
 Luna\dots\ }\end{center}
\switchcolumn
\selectlanguage{english}
\begin{center}{\color{gregoriocolor} The 29\textsuperscript{th} Day of 
 January. The\dots\ Day of the Moon.}\end{center}
\end{paracol}

\noindent\begin{tabularx}{\linewidth}{*{19}{>{\centering\arraybackslash}X}}
 \textcolor{gregoriocolor}{a} & \textcolor{gregoriocolor}{b} & \textcolor{gregoriocolor}{c} & \textcolor{gregoriocolor}{d} & \textcolor{gregoriocolor}{e} & \textcolor{gregoriocolor}{f} & \textcolor{gregoriocolor}{g} & \textcolor{gregoriocolor}{h} & \textcolor{gregoriocolor}{i} & \textcolor{gregoriocolor}{k} & \textcolor{gregoriocolor}{l} & \textcolor{gregoriocolor}{m} & \textcolor{gregoriocolor}{n} & \textcolor{gregoriocolor}{p} & \textcolor{gregoriocolor}{q} & \textcolor{gregoriocolor}{r} & \textcolor{gregoriocolor}{s} & \textcolor{gregoriocolor}{t} & \textcolor{gregoriocolor}{u} \\
 30 & 1 & 2 & 3 & 4 & 5 & 6 & 7 & 8 & 9 & 10 & 11 & 12 & 13 & 14 & 15 & 16 & 17 & 18 \\
\end{tabularx}
\vspace{0.5\baselineskip}
\noindent\begin{tabularx}{\linewidth}{*{12}{>{\centering\arraybackslash}X}}
 \textcolor{gregoriocolor}{A} & \textcolor{gregoriocolor}{B} & \textcolor{gregoriocolor}{C} & \textcolor{gregoriocolor}{D} & \textcolor{gregoriocolor}{E} & F & \textcolor{gregoriocolor}{F} & \textcolor{gregoriocolor}{G} & \textcolor{gregoriocolor}{H} & \textcolor{gregoriocolor}{M} & \textcolor{gregoriocolor}{N} & \textcolor{gregoriocolor}{P} \\
 19 & 20 & 21 & 22 & 23 & 24 & 24 & 25 & 26 & 27 & 28 & 29 \\
\end{tabularx}

\begin{paracol}{2}
\selectlanguage{latin}
\lettrine[lines=2]{S}{ancti} Francísci Salésii, Epíscopi Gebennénsis, Confessóris et Ecclésiæ 
 Doctóris, ómnium Scriptórum catholicórum, diáriis aliísve scriptis in vulgus 
 edéndis sapiéntiam Christiánam illustrántium ac provehéntium et tutántium, 
 peculiáris apud Deum Patróni; qui migrávit in cælum quinto Kaléndas Januárii, 
 sed hac die, ob Translatiónem córporis ejus, potíssimum cólitur.
\switchcolumn
\selectlanguage{english}
\lettrine[lines=2]{S}{t.} Francis de Sales, bishop of Geneva, confessor and doctor of the Church, 
 special patron before God of all Catholic writers in explaining, promoting, 
 or defending Christian doctrine either by publishing journals or other 
 writings in the vernacular. He departed to heaven on the 28th of 
 December, but because of the transfer of his body on this day, his feast is 
 now celebrated.
\switchcolumn*
\selectlanguage{latin}
Tréviris deposítio beáti Valérii Epíscopi, qui fuit discípulus sancti Petri 
 Apóstoli.
\switchcolumn
\selectlanguage{english}
At Treves, the death of the blessed bishop Valerius, disciple of the apostle 
 St. Peter.
\switchcolumn*
\selectlanguage{latin}
Romæ, via Nomentána, natális sanctórum Mártyrum 
 Pápiæ et Mauri mílitum, témpore Diocletiáni Imperatóris; quorum ora jussit 
 Laodícius, Urbis Præféctus, ad primam Christi confessiónem lapídibus 
 contúndi, et sic eos in cárcerem trahi, ac póstea fústibus cædi, atque ad 
 últimum plumbátis pércuti, donec exspirárent.
\switchcolumn
\selectlanguage{english}
At Rome, on the Via Nomentana, the birthday of the holy martyrs Papias and 
 Maur, soldiers under Emperor Diocletian. At their first confession of 
 Christ they had their mouths bruised with stones and were thrown into prison 
 by order of Laodicius, prefect of the city. Afterwards they were 
 beaten with rods and with leaded whips until they expired.
\switchcolumn*
\selectlanguage{latin}
Perúsiæ sancti Constántii, Epíscopi et Mártyris; 
 qui, una cum Sóciis, sub Marco Aurélio Imperatóre, ob fídei defensiónem, 
 martyrii corónam accépit.
\switchcolumn
\selectlanguage{english}
At Perugia, in the time of Marcus Aurelius, St. Constantius, bishop and 
 martyr, who, together with his companions, received the crown of martyrdom 
 for the defence of the faith.
\switchcolumn*
\selectlanguage{latin}
Medioláni sancti Aquilíni Presbyteri, qui, ab Ariánis gládio in gútture 
 transfíxus, martyrio coronátur.
\switchcolumn
\selectlanguage{english}
At Milan, St. Aquilinus, priest, who was crowned with martyrdom by having 
 his throat pierced with a sword by the Arians.
\switchcolumn*
\selectlanguage{latin}
Edéssæ, in Syria, sanctórum Mártyrum Sarbélii 
 et Bárbeæ soróris, qui, a beáto Barsimǽo Epíscopo baptizáti, ambo, in 
 persecutióne Trajáni, sub Lysia Præside, martyrio coronáti sunt.
\switchcolumn
\selectlanguage{english}
At Edessa in Syria, the holy martyrs Sabellus and his sister Barbea, who 
 were baptized by the blessed bishop Barsimaeus, and crowned with martyrdom 
 in the persecution of Trajan, under the governor Lysias.
\switchcolumn*
\selectlanguage{latin}
In território Tricassíno sancti Sabiniáni Mártyris, qui, jubénte Aureliáno 
 Imperatóre, pro fide Christi decollátus est.
\switchcolumn
\selectlanguage{english}
In the territory of Troyes, St. Sabinian, martyr, who was beheaded for the 
 faith of Christ by command of the emperor Aurelian.
\switchcolumn*
\selectlanguage{latin}
Apud Bitúricas, in Aquitánia, sancti Sulpícii Sevéri 
 Epíscopi, virtútibus et eruditióne conspícui.
\switchcolumn
\selectlanguage{english}
At Bourges, St. Sulpice Severus, bishop, distinguished by his virtues and 
 learning.
\switchcolumn*
\selectlanguage{latin}
\end{paracol}


% ---- martyrology/mart01/mart0130.htm
\needspace{10\baselineskip}
\begin{paracol}{2}
\selectlanguage{latin}
\begin{center}{\color{gregoriocolor} Tértio Kaléndas Februárii. 
 Luna\dots\ }\end{center}
\switchcolumn
\selectlanguage{english}
\begin{center}{\color{gregoriocolor} The 30\textsuperscript{th} Day of 
 January. The\dots\ Day of the Moon.}\end{center}
\end{paracol}

\noindent\begin{tabularx}{\linewidth}{*{19}{>{\centering\arraybackslash}X}}
 \textcolor{gregoriocolor}{a} & \textcolor{gregoriocolor}{b} & \textcolor{gregoriocolor}{c} & \textcolor{gregoriocolor}{d} & \textcolor{gregoriocolor}{e} & \textcolor{gregoriocolor}{f} & \textcolor{gregoriocolor}{g} & \textcolor{gregoriocolor}{h} & \textcolor{gregoriocolor}{i} & \textcolor{gregoriocolor}{k} & \textcolor{gregoriocolor}{l} & \textcolor{gregoriocolor}{m} & \textcolor{gregoriocolor}{n} & \textcolor{gregoriocolor}{p} & \textcolor{gregoriocolor}{q} & \textcolor{gregoriocolor}{r} & \textcolor{gregoriocolor}{s} & \textcolor{gregoriocolor}{t} & \textcolor{gregoriocolor}{u} \\
 1 & 2 & 3 & 4 & 5 & 6 & 7 & 8 & 9 & 10 & 11 & 12 & 13 & 14 & 15 & 16 & 17 & 18 & 19 \\
\end{tabularx}
\vspace{0.5\baselineskip}
\noindent\begin{tabularx}{\linewidth}{*{12}{>{\centering\arraybackslash}X}}
 \textcolor{gregoriocolor}{A} & \textcolor{gregoriocolor}{B} & \textcolor{gregoriocolor}{C} & \textcolor{gregoriocolor}{D} & \textcolor{gregoriocolor}{E} & F & \textcolor{gregoriocolor}{F} & \textcolor{gregoriocolor}{G} & \textcolor{gregoriocolor}{H} & \textcolor{gregoriocolor}{M} & \textcolor{gregoriocolor}{N} & \textcolor{gregoriocolor}{P} \\
 20 & 21 & 22 & 23 & 24 & 25 & 25 & 26 & 27 & 28 & 29 & 30 \\
\end{tabularx}

\begin{paracol}{2}
\selectlanguage{latin}
\lettrine[lines=2]{S}{anctæ} Martínæ, Vírginis et Mártyris, cujus 
 dies natális Kaléndis Januárii recólitur.
\switchcolumn
\selectlanguage{english}
\lettrine[lines=2]{S}{t.} Martina, virgin and martyr, who is commemorated on her birthday, the 
 first day of this month.
\switchcolumn*
\selectlanguage{latin}
Edéssæ, in Syria, sancti Barsimǽi Epíscopi, 
 qui, cum Gentíles plúrimos convertísset ad fidem et præmisísset ad corónam, 
 eos secútus est, sub Trajáno, cum palma martyrii.
\switchcolumn
\selectlanguage{english}
At Edessa in Syria, in the reign of Trajan, St. Barsimaeus, bishop, who 
 converted many Gentiles to the faith, sent them before him to gain their 
 crown, and then followed them with the palm of martyrdom.
\switchcolumn*
\selectlanguage{latin}
Antiochíæ pássio beáti Hippólyti Presbyteri, 
 qui, decéptus aliquándiu schísmate Nováti, sed, operánte grátia Christi, 
 corréctus, ad unitátem Ecclésiæ rédiit, pro qua et in qua póstea illústre 
 martyrium consummávit. Hic, rogátus a suis quænam secta vérior esse 
 servándam quam Petri Cáthedra custodíret, júgulum præbuit.
\switchcolumn
\selectlanguage{english}
At Antioch, the passion of the blessed Hippolytus, priest, who for a short 
 time deceived by the Novatian schism, was converted by the grace of Christ, 
 and returned to the unity of the Church, for which and in which he 
 afterwards underwent a glorious martyrdom. Being asked by the 
 schismatics, which was the better side, he said that he detested the 
 doctrine of Novatus, and that the faith which the Chair of Peter taught 
 ought to be professed, after which he was beheaded.
\switchcolumn*
\selectlanguage{latin}
In Africa pássio sanctórum Mártyrum Feliciáni, Philappiáni et aliórum centum 
 vigínti quátuor.
\switchcolumn
\selectlanguage{english}
In Africa, the passion of the holy martyrs Felician, Philappian, and one 
 hundred and twenty-four others.
\switchcolumn*
\selectlanguage{latin}
Item beáti Alexándri, qui, in persecutióne Décii, comprehénsus est, ac, longævæ 
 ætátis veneránda canítie et confessióne iteráta respléndens, inter 
 carníficum torménta réddidit spíritum.
\switchcolumn
\selectlanguage{english}
Blessed Alexander, a man of venerable aspect and advanced age, who was 
 apprehended in the persecution of Decius. After gloriously and 
 repeatedly confessing the faith, in the midst of torments he gave up his 
 soul unto God.
\switchcolumn*
\selectlanguage{latin}
Edéssæ, in Syria, sancti Barsis Epíscopi, dono 
 curatiónum illústris; qui, a Valénte, Imperatóre Ariáno, in díssitas 
 regiónes ob fidem cathólicam relegátus, ac tríplici mutatióne fatigátus 
 exsílii, vitam finívit.
\switchcolumn
\selectlanguage{english}
At Edessa in Syria, St Barses, bishop, renowned for the gift of healing 
 diseases. For holding to the Catholic faith he was banished by the 
 Arian emperor Valens into the most remote corner of that country, and he 
 there ended his days.
\switchcolumn*
\selectlanguage{latin}
Hierosólymis natális sancti Matthíæ Epíscopi, 
 de quo mira et plena fídei gesta narrántur; qui, sub Hadriáno, multa pro 
 Christo perpéssus est, ac demum in pace quiévit.
\switchcolumn
\selectlanguage{english}
At Jerusalem, the birthday of St. Matthias, bishop, of whom wonderful deeds 
 are related which were inspired by faith. After having endured many 
 trials for Christ under Adrian, he passed away in peace.
\switchcolumn*
\selectlanguage{latin}
Papíæ sancti Armentárii, Epíscopi et 
 Confessóris.
\switchcolumn
\selectlanguage{english}
At Pavia, St. Armentarius, bishop and confessor.
\switchcolumn*
\selectlanguage{latin}
In Malbódio, Hannóniæ monastério, sanctæ 
 Aldegúndis Vírginis, témpore Dagobérti Regis.
\switchcolumn
\selectlanguage{english}
In Hainaut, in the monastery of Maubeuge, St. Aldegund, virgin, who lived in 
 the time of King Dagobert.
\switchcolumn*
\selectlanguage{latin}
Vitérbii sanctæ Hyacínthæ de Mariscóttis 
 Vírginis, ex tértio sancti Francísci Ordine Sanctimoniális, pæniténtia et 
 caritáte insígnis; quam Pius Papa Séptimus Sanctis adscrípsit.
\switchcolumn
\selectlanguage{english}
At Viterbo, the holy virgin Hyacinth Mariscotti, a nun of the Third Order of 
 St. Francis, distinguished for the virtues of penance and charity. She 
 was inscribed among the saints by Pope Pius VII.
\switchcolumn*
\selectlanguage{latin}
Medioláni sanctæ Savínæ, féminæ religiosíssimæ, 
 quæ, ad sepúlcra sanctórum Náboris et Felícis Mártyrum orans, obdormívit in 
 Dómino.
\switchcolumn
\selectlanguage{english}
At Milan, St. Savina, a most religious woman, who went to rest in the Lord 
 while praying at the tomb of the holy martyrs Nabor and Felix.
\switchcolumn*
\selectlanguage{latin}
In território Parisiénsi sanctæ Bathíldis 
 Regínæ, sanctitáte et miraculórum glória præcláræ.
\switchcolumn
\selectlanguage{english}
In the district of Paris, St. Bathilde, queen, renowned for the worthiness 
 of her miracles and her sanctity.
\switchcolumn*
\selectlanguage{latin}
\end{paracol}


% ---- martyrology/mart01/mart0131.htm
\needspace{10\baselineskip}
\begin{paracol}{2}
\selectlanguage{latin}
\begin{center}{\color{gregoriocolor} Prídie Kaléndas Februárii. 
 Luna\dots\ }\end{center}
\switchcolumn
\selectlanguage{english}
\begin{center}{\color{gregoriocolor} The 31\textsuperscript{st} Day of 
 January. The\dots\ Day of the Moon.}\end{center}
\end{paracol}

\noindent\begin{tabularx}{\linewidth}{*{19}{>{\centering\arraybackslash}X}}
 \textcolor{gregoriocolor}{a} & \textcolor{gregoriocolor}{b} & \textcolor{gregoriocolor}{c} & \textcolor{gregoriocolor}{d} & \textcolor{gregoriocolor}{e} & \textcolor{gregoriocolor}{f} & \textcolor{gregoriocolor}{g} & \textcolor{gregoriocolor}{h} & \textcolor{gregoriocolor}{i} & \textcolor{gregoriocolor}{k} & \textcolor{gregoriocolor}{l} & \textcolor{gregoriocolor}{m} & \textcolor{gregoriocolor}{n} & \textcolor{gregoriocolor}{p} & \textcolor{gregoriocolor}{q} & \textcolor{gregoriocolor}{r} & \textcolor{gregoriocolor}{s} & \textcolor{gregoriocolor}{t} & \textcolor{gregoriocolor}{u} \\
 2 & 3 & 4 & 5 & 6 & 7 & 8 & 9 & 10 & 11 & 12 & 13 & 14 & 15 & 16 & 17 & 18 & 19 & 20 \\
\end{tabularx}
\vspace{0.5\baselineskip}
\noindent\begin{tabularx}{\linewidth}{*{12}{>{\centering\arraybackslash}X}}
 \textcolor{gregoriocolor}{A} & \textcolor{gregoriocolor}{B} & \textcolor{gregoriocolor}{C} & \textcolor{gregoriocolor}{D} & \textcolor{gregoriocolor}{E} & F & \textcolor{gregoriocolor}{F} & \textcolor{gregoriocolor}{G} & \textcolor{gregoriocolor}{H} & \textcolor{gregoriocolor}{M} & \textcolor{gregoriocolor}{N} & \textcolor{gregoriocolor}{P} \\
 21 & 22 & 23 & 24 & 25 & 26 & 26 & 27 & 28 & 29 & 30 & 1 \\
\end{tabularx}

\begin{paracol}{2}
\selectlanguage{latin}
\lettrine[lines=2]{A}{ugústæ} Taurinórum sancti Joánnis Bosco, 
 Confessóris, Societátis Salesiánæ ac Institúti Filiárum Maríæ Auxiliatrícis 
 Fundatóris, animárum zelo et fídei propagándæ conspícui, quem Pius Papa 
 Undécimus Sanctórum fastis adscrípsit.
\switchcolumn
\selectlanguage{english}
\lettrine[lines=2]{A}{t} Turin,the birthday of St. John Bosco, confessor, founder of the Salesian 
 Congregation and of the Institute of the Daughters of Mary, Help of 
 Christians. Conspicuous for his zeal for souls and for the propagation 
 of the faith, he was canonized by Pope Pius XI.
\switchcolumn*
\selectlanguage{latin}
Romæ, via Portuénsi, sanctórum Mártyrum Cyri et 
 Joánnis, qui pro confessióne Christi, post multa torménta, cápite truncáti 
 sunt.
\switchcolumn
\selectlanguage{english}
At Rome, on the road to Ostia, the holy martyrs Cyrus and John, who were 
 beheaded after suffering many torments for the name of Christ.
\switchcolumn*
\selectlanguage{latin}
Alexandríæ natális sancti Metráni Mártyris, 
 qui, sub Décio Imperatóre, cum ad jussiónem Paganórum nollet ímpia verba 
 proférre, hi totum ejus corpus fústibus collisérunt, vultúmque et óculos 
 præacútis cálamis terebrántes, cum cruciátibus expulérunt ipsum extra urbem, 
 ibíque lapídibus oppréssum interemérunt.
\switchcolumn
\selectlanguage{english}
At Alexandria, in the time of Emperor Decius, the birthday of St. Metran, 
 martyr, who, because he refused to utter blasphemous words at the bidding of 
 the pagans, had his body all bruised with blows, and his face and eyes 
 pierced with sharp pointed reeds. He was then driven out of the city 
 and stoned to death.
\switchcolumn*
\selectlanguage{latin}
Ibídem sanctórum Mártyrum Saturníni, Thyrsi et Victóris.
\switchcolumn
\selectlanguage{english}
In the same place, the holy martyrs Saturninus, Thyrsus, and Victor.
\switchcolumn*
\selectlanguage{latin}
Item Alexandríæ sanctórum Mártyrum Tharsícii, 
 Zótici, Cyríaci et Sociórum.
\switchcolumn
\selectlanguage{english}
Also at Alexandria, the holy martyrs Tharsicius, Zoticus, Cyriacus, and 
 their companions.
\switchcolumn*
\selectlanguage{latin}
Cyzici, in Hellespónto, sanctæ Tryphǽnæ 
 Mártyris, quæ, plúrimis torméntis superátis, a tauro demum necáta, martyrii 
 palmam proméruit.
\switchcolumn
\selectlanguage{english}
At Cyzicum in the Hellespont, St. Triphenes, martyr, who overcame various 
 torments, but was finally killed by a bull, and thus merited the palm of 
 martyrdom.
\switchcolumn*
\selectlanguage{latin}
Mútinæ sancti Geminiáni Epíscopi, miraculórum 
 glória conspícui.
\switchcolumn
\selectlanguage{english}
At Modena, St. Geminian, bishop, made illustrious by his miracles.
\switchcolumn*
\selectlanguage{latin}
In Província Mediolanénsi sancti Júlii, Presbyteri et Confessóris, témpore 
 Imperatóris Theodósii.
\switchcolumn
\selectlanguage{english}
In the province of Milan, St. Julius, priest and confessor, in the reign of 
 the emperor Theodosius.
\switchcolumn*
\selectlanguage{latin}
Neápoli sancti Francísci Xavérii-Maríæ Biánchi, 
 Confessóris, Clérici reguláris sancti Pauli, signis, donis cæléstibus et 
 admirábili patiéntia illústris, quem Pius Papa Duodécimus ad suprémos 
 honóres Sanctórum éxtulit.
\switchcolumn
\selectlanguage{english}
At Naples, St. Francis Xavier-Maria Bianchi, confessor, cleric regular of 
 St. Paul, renowned for miracles, heavenly gifts and an admirable patience, 
 whom Pope Pius XII raised to the supreme honour of sainthood.
\switchcolumn*
\selectlanguage{latin}
Romæ sanctæ Marcéllæ Víduæ, cujus præcláras 
 laudes beátus Hierónymus scripsit.
\switchcolumn
\selectlanguage{english}
At Rome, St. Marcella, widow, whose meritorious deeds are related by St. 
 Jerome.
\switchcolumn*
\selectlanguage{latin}
Item Romæ beátæ Ludovícæ Albertóniæ, Víduæ 
 Románæ, ex tértio Ordine sancti Francísci, virtútibus claræ.
\switchcolumn
\selectlanguage{english}
Also at Rome, blessed Louise Albertonia, a Roman widow, member of the Third 
 Order of St. Francis, distinguished for her virtues.
\switchcolumn*
\selectlanguage{latin}
Eódem die Translátio sancti Marci Evangelístæ, 
 cum sacrum ejus corpus ex Alexandría, a bárbaris tunc occupáta, Venétias 
 allátum, ibídem in majóri Ecclésia, ejus nómine consecráta, 
 honorificentíssime cónditum fuit.
\switchcolumn
\selectlanguage{english}
The same day, the transfer of the revered body of the Evangelist St. Mark 
 from the city of Alexandria in Egypt, then occupied by barbarians, to 
 Venice, and with the greatest honours placed in the large church dedicated 
 to his name.
\switchcolumn*
\selectlanguage{latin}
\end{paracol}
